\section{非利士-北方军队+神降火=荒凉毀滅，成为犹大的牧场 OK}
\begin{figure}[H] 
\centering
\includegraphics[width=6.in]{ZongJie/figs/JiaSha}
\end{figure}
\subsection{以賽亞書 14:28-32-神以饑荒治死你的根，餘剩的人必被北方出來的殺戮}
亞哈斯王崩的那年，就有以下的默示。
非利士全地啊，不要因擊打你的杖折斷就喜樂，因為從蛇的根必生出毒蛇，牠所生的是火焰的飛龍。 30 貧寒人的長子必有所食，窮乏人必安然躺臥。我必以饑荒治死你的根，你所餘剩的人必被殺戮。 31 門哪，應當哀號！城啊，應當呼喊！非利士全地啊，你都消化了！因為有煙從北方出來，他行伍中必無亂隊的。」
32 可怎樣回答外邦 (非利士)的使者呢？必說：「耶和華建立了錫安，他百姓中的困苦人必投奔在其中。」

\subsection{耶利米書 47:1-22-北方漲溢的河毀滅非利士人}
法老攻擊加沙之先，有耶和華論非利士人的話臨到先知耶利米。2 耶和華如此說：「有水從北方發起，成為漲溢的河，要漲過遍地和其中所有的，並城和其中所住的。人必呼喊，境內的居民都必哀號。 3 聽見敵人壯馬蹄跳的響聲和戰車隆隆、車輪轟轟，為父的手就發軟，不回頭看顧兒女。 4 因為日子將到，要\textcolor{red}{毀滅一切非利士人}，剪除幫助推羅、西頓所剩下的人。原來耶和華必\textcolor{red}{毀滅非利士人}，就是迦斐托海島餘剩的人。 5 加沙成了光禿，平原中所剩的亞實基倫歸於無有。你用刀劃身，要到幾時呢？ 6 耶和華的刀劍哪，你到幾時才止息呢？你要入鞘，安靖不動！ 7 耶和華既吩咐你攻擊亞實基倫和海邊之地，他已經派定你，焉能止息呢？」

\subsection{以西結書 25:15-17-主耶和華必伸手攻擊非利士滅絕沿海剩下的居民}
主耶和華如此說：因非利士人向猶大人報仇，就是以恨惡的心報仇雪恨，永懷仇恨，要毀滅他們， 16 所以主耶和華如此說：我必伸手\textcolor{red}{攻擊非利士人}，剪除基利提人，滅絕沿海剩下的居民。 17 我向他們大施報應，發怒斥責他們。我報復他們的時候，他們就知道我是耶和華。

\subsection{阿摩司書 1:6-8-降火在加沙非利士人所餘剩的必都滅亡}
耶和華如此說：「加沙三番四次地犯罪，我必不免去他的刑罰，因為他擄掠眾民交給以東。 7 我卻要\textcolor{red}{降火在加沙}的城內，燒滅其中的宮殿。 8 我必\textcolor{red}{剪除亞實突的居民}和\textcolor{red}{亞實基倫掌權}的，也必\textcolor{red}{反手攻擊以革倫}，非利士人所餘剩的必都\textcolor{red}{滅亡}。」這是主耶和華說的。

\subsection{西番雅書 2:4-7-加沙必致見棄，沿海之地要變為猶大家剩下人的草場}
加沙必致見棄，亞實基倫必然荒涼。人在正午必趕出亞實突的民，以革倫也被拔出根來。 5 住沿海之地的基利提族有禍了！迦南，非利士人之地啊，耶和華的話與你反對，說：「我必毀滅你，以致無人居住。」 6 沿海之地要變為草場，其上有牧人的住處和羊群的圈。 7 這地必為猶大家剩下的人所得，他們必在那裡牧放群羊，晚上必躺臥在亞實基倫的房屋中，因為耶和華他們的神必眷顧他們，使他們被擄的人歸回。

\subsection{撒迦利亞書 9:1-7-加沙必不再有君王，除滅非利士人的驕傲}
耶和華的默示應驗在哈得拉地大馬士革——世人和以色列各支派的眼目都仰望耶和華—— 2 和靠近的哈馬，並推羅、西頓，因為這二城的人大有智慧。 3 「推羅為自己修築保障，積蓄銀子如塵沙，堆起精金如街上的泥土。 4 主必趕出她，打敗她海上的權力，她必被火燒滅。 
5 亞實基倫看見必懼怕，加沙看見甚痛苦，以革倫因失了盼望蒙羞。加沙必不再有君王，亞實基倫也不再有居民。 6 私生子(外族人)必住在亞實突，我必除滅非利士人的驕傲。 7 我必除去他口中帶血之肉和牙齒內可憎之物，他必作為餘剩的人歸於我們的神，必在猶大像族長，以革倫人必如耶布斯人。








\section{亞捫}
\subsection{亞捫背景}
\subsubsection{亞捫来源}
2.亚扪－是”先天”的意思，其原名便雅米(意”我民之子”)，就是以后亚扪的始祖。到了以色列人进迦南地的时候，成为大族，其人民最残忍(撒上11:2，摩1:13)。他们住在约但河的东边，南境与北境连接拉巴城(撒下11:1，结25:5)。以色列人逼近迦南时，亚扪人拒绝他们经过其属地，且请巴兰咒诅以色列人，所以神永远不准他们入耶和华的会。所罗门王时，纳亚扪族的女儿拿玛为妃，生下罗波安(王上14:21)。亚扪人所敬拜的神名叫米勒公，又叫摩洛，因所罗门王娶了摩押和亚扪女子为妃嫔，而使他陷于敬拜这些偶像的罪中(王上11:1，4－8)。亚扪人的主城是拉巴，或者被称作“拉巴亚扪“（意为亚扪人的都市）。后来希腊人将这座城市称为“拉巴安曼那”。当代约旦国的首都安曼大致位于古代拉巴城址之上。圣经预言摩押和亚扪时，却提到这两国还要存在，两国有人要归回，“耶和华说，到末后，我还要使被掳的摩押人归回。”“后来我还要使被掳的亚扪人归回。这是耶和华说的”(耶48:47；49:6)。

\subsubsection{亞捫人的首都 - 拉巴Rabbah,(約旦的首都阿曼 Amman)}

有時也稱為「拉巴亞捫」。位於約但河東岸、希實本東北約 21 公里，雅博河的發源處，水泉甚多，故有水城之稱。土壤肥沃，宜於農牧，一直都是軍事和商業重鎮，也一直是亞捫人的首都。該城在古時即有君王大道經過，交通便利.靠近亞羅珥。

以色列人擊敗巴珊王噩，他的鐵床(有說是棺槨)留在拉巴。以色列人分配地業時，劃歸迦得支派，但亞捫人仍佔據該地。

大衛時期，他的將軍約押才攻取了拉巴；赫人烏利亞即因謀殺而死於此役(因大衛看中其妻拔示巴)。大衛逃避押沙龍時，曾得拉巴人的幫助。先知們斥責拉巴的惡行，並預言它的毀滅。在大衛王時代，城區分為上下兩城，上城即是京城，下城即是水城。

從 1924 至 1947 年間的考古工作得知，該地在舊石器時代，就是在一萬年以前已有人類居住，遺物中掘出有主前八千年代中所建的石架和巨石紀念碑，城牆約是在 2300 BC 建成，另有喜克索時代(1600 BC)的水池，1400 BC 前後所建之神廟，還有許多不同時代的墳墓。因為現今之城是建造在舊城之上，除了在地面上尚存有的希臘羅馬時代劇場和神殿等的遺物之外，使得必須靠挖掘的考古工作不易進行，其更詳細的歷史無法作進一步的考證，只有從其他之資料來了解。主前三世紀，拉巴成為羅馬帝國的屬地，埃及的多利買二世，重建此城，改名叫非拉鐵非，在新約時代是十城之一。阿拉伯人攻佔領拉巴後，此地逐漸衰微。

\begin{table}[H]
\centering 
\begin{tabular}{p{2.cm}|p{15.5cm}}\hline\hline
申3:11&利乏音人所剩下的，只有巴珊王噩，他的床是鐵的，長九肘，寬四肘，現今不是在亞捫人的拉巴麼。\\\hline

書13:24-25&摩西按著迦得支派的宗族，分給他們產業，他們的境界是雅謝和基列各城，並亞捫人一半的地，直到拉巴前的亞羅珥。\\\hline

撒下11:1-17&\\\hline
代上20:1-3&過了一年，到列王出戰的時候，大衛又差派約押率領以色列人出戰，他們就打敗亞捫人，毀壞亞捫人的地，圍攻拉巴。\\\hline

撒下12:26-31&約押攻打亞捫人的京城拉巴，取其水城，大衛王再領兵來，就奪取了拉巴。\\\hline

撒下17:27&大衛因逃避押沙龍到了瑪哈念，亞捫族的拉巴人拿轄的兒子朔比等，帶了用具和食物來給大衛和他的隨從。\\\hline

耶49:2&先知預言亞捫必受罰，拉巴必遭攻擊，拉巴要成為亂堆。\\\hline

結21:20&先知預言巴比倫王必攻擊拉巴及耶路撒冷。\\\hline

結25:5&先知的預言說，我必使拉巴為駱駝場，使亞捫人的地成為羊群躺臥之處。\\\hline

摩1:14&先知的預言說，我卻要在爭戰吶喊的日子，旋風狂暴的時候，點火在拉巴的城內，燒滅其中的宮殿。\\\hline
\end{tabular}
\caption{拉巴}
\end{table}


\subsection{耶利米書-預言亞捫人必受懲罰後必復興 49:1-6}
%\subsubsection{預言亞捫人必受懲罰後必復興 49:1-6}
論亞捫人。耶和華如此說：「以色列沒有兒子嗎？沒有後嗣嗎？瑪勒堪為何得迦得之地為業呢？屬他的民為何住其中的城邑呢？」 2 耶和華說：「日子將到，我必使人聽見打仗的喊聲，是攻擊亞捫人拉巴的喊聲。拉巴要成為亂堆，屬她的鄉村(女子)要被火焚燒。先前得以色列地為業的，此時以色列倒要得他們的地為業。」這是耶和華說的。 3 「希實本哪，你要哀號！因為愛地變為荒場。拉巴的居民 (女子)哪，要呼喊！以麻布束腰，要哭號，在籬笆中跑來跑去，因瑪勒堪和屬他的祭司、首領要一同被擄去。 4 背道的民(女子)哪，你們為何因有山谷，就是水流的山谷誇張呢？為何倚靠財寶說『誰能來到我們這裡』呢？」 5 主萬軍之耶和華說：「我要使恐嚇從四圍的人中臨到你們，你們必被趕出，各人一直前往，沒有人收聚逃民。「後來，我還要使被擄的亞捫人歸回。」這是耶和華說的

\subsection{以西結書}
\subsubsection{預言亞捫人亦將遭報 21:28-32}
「人子啊，要發預言說：『主耶和華論到亞捫人和他們的凌辱，吩咐我如此說：有刀，有拔出來的刀，已經擦亮，為行殺戮，使它像閃電以行吞滅。 29 人為你見虛假的異象，行謊詐的占卜，使你倒在受死傷之惡人的頸項上。他們罪孽到了盡頭，受報的日子已到。 30 你將刀收入鞘吧！在你受造之處，生長之地，我必刑罰你。 31 我必將我的惱恨倒在你身上，將我烈怒的火噴在你身上，又將你交在善於殺滅的畜類人手中。 32 你必當柴被火焚燒，你的血必流在國中，你必不再被記念，因為這是我耶和華說的。』」

\subsubsection{預言亞捫人樂以色列遘災必受重報 25:1-15}
耶和華的話臨到我說： 2 「人子啊，你要面向亞捫人說預言攻擊他們， 3 說：『你們當聽主耶和華的話！主耶和華如此說：我的聖所被褻瀆，以色列地變荒涼，猶大家被擄掠，那時你便因這些事說「啊哈！」。 4 所以我必將你的地交給東方人為業，他們必在你的地上安營居住，吃你的果子，喝你的奶。 5 我必使拉巴為駱駝場，使亞捫人的地為羊群躺臥之處，你們就知道我是耶和華。 6 主耶和華如此說：因你拍手頓足，以滿心的恨惡向以色列地歡喜， 7 所以我伸手攻擊你，將你交給列國作為擄物。我必從萬民中剪除你，使你從萬國中敗亡。我必除滅你，你就知道我是耶和華。』

%\subsubsection{摩押，以東，亞捫亦必遭報 25:8-15}
主耶和華如此說：因摩押和西珥人說『看哪，猶大家與列國無異』， 9 所以我要破開摩押邊界上的城邑，就是摩押人看為本國之榮耀的伯耶西末、巴力免、基列亭， 10 好使東方人來攻擊亞捫人。我必將亞捫人之地交給他們為業，使亞捫人在列國中不再被記念。 11 我必向摩押施行審判，他們就知道我是耶和華。

\subsection{阿摩司書 1:13-15}
耶和華如此說：「亞捫人三番四次的犯罪，我必不免去他們的刑罰，因為他們剖開基列的孕婦，擴張自己的境界。 14 我卻要在爭戰呐喊的日子，旋風狂暴的時候，點火在拉巴的城內，燒滅其中的宮殿。 15 他們的王和首領必一同被擄去。」這是耶和華說的。

\subsection{西番雅書 2:8-11}
我聽見摩押人的毀謗和亞捫人的辱罵，就是毀謗我的百姓，自誇自大，侵犯他們的境界。」 9 萬軍之耶和華以色列的神說：「我指著我的永生起誓：摩押必像所多瑪，亞捫人必像蛾摩拉，都變為刺草、鹽坑、永遠荒廢之地。我百姓所剩下的必擄掠他們，我國中所餘剩的必得著他們的地。」 10 這事臨到他們，是因他們驕傲，自誇自大，毀謗萬軍之耶和華的百姓。 11 耶和華必向他們顯可畏之威，因他必叫世上的諸神瘦弱，列國海島的居民各在自己的地方敬拜他。










\section{对亚兰之审判}
\subsection{以賽亞書 17:1-3}
論大馬士革的默示。

「看哪，大馬士革已被廢棄，不再為城，必變做亂堆。 2 亞羅珥的城邑已被撇棄，必成為牧羊之處，羊在那裡躺臥，無人驚嚇。 3 以法蓮不再有保障，大馬士革不再有國權，亞蘭所剩下的必像以色列人的榮耀消滅一樣。」這是萬軍之耶和華說的。

\subsection{耶利米書-大馬士革必受懲罰 49:23-27}
%\subsubsection{大馬士革必受懲罰 49:1-22}
論大馬士革。「哈馬和亞珥拔蒙羞，因他們聽見凶惡的信息就消化了。海上有憂愁，不得平靜。 24 大馬士革發軟，轉身逃跑。戰兢將她捉住，痛苦憂愁將她抓住如產難的婦人一樣。 25 我所喜樂、可稱讚的城，為何被撇棄了呢？ 26 她的少年人必仆倒在街上，當那日一切兵丁必默默無聲。」這是萬軍之耶和華說的。 27 「我必在大馬士革城中使火著起，燒滅便哈達的宮殿。」

\subsection{阿摩司書 1:3-5}
\subsubsection{1:3-5}
耶和華如此說：「大馬士革三番四次地犯罪，我必不免去他的刑罰，因為他以打糧食的鐵器打過基列。 4 我卻要降火在哈薛的家中，燒滅便哈達的宮殿。 5 我必折斷大馬士革的門閂，剪除亞文平原的居民和伯伊甸掌權的，亞蘭人必被擄到吉珥。」這是耶和華說的。

\subsubsection{擄到大馬士革以外 5:25-27}
以色列家啊，你們在曠野四十年，豈是將祭物和供物獻給我呢？ 26 你們抬著為自己所造之摩洛的帳幕和偶像的龕，並你們的神星。 27 所以，我要把你們擄到大馬士革以外。」這是耶和華，名為萬軍之神說的。




\subsection{撒迦利亞書 9:1-7}
耶和華的默示應驗在哈得拉地大馬士革——世人和以色列各支派的眼目都仰望耶和華—— 2 和靠近的哈馬，並推羅、西頓，因為這二城的人大有智慧。 3 「推羅為自己修築保障，積蓄銀子如塵沙，堆起精金如街上的泥土。 4 主必趕出她，打敗她海上的權力，她必被火燒滅。 5 亞實基倫看見必懼怕，加沙看見甚痛苦，以革倫因失了盼望蒙羞。加沙必不再有君王，亞實基倫也不再有居民。 6 私生子(外族人)必住在亞實突，我必除滅非利士人的驕傲。 7 我必除去他口中帶血之肉和牙齒內可憎之物，他必作為餘剩的人歸於我們的神，必在猶大像族長，以革倫人必如耶布斯人。







\section{論古實之預言}
\subsection{古實-以斯拉培训网路$（http://www.equiptoserve.org/）$}
埃及南方的古實。希伯來文的 「古實」由古埃及語借來，古實原為人名，他是挪亞的孫子、含的兒子。古實的後代稱為「古實人」，他們居住的地方被稱為「古實地」。《舊約》中有四十二處提及古實、古實人、古實河；新約時代按希臘語發音，稱 「埃提阿伯」(今譯埃塞俄比亞)。該地在非洲東部、尼羅河上游東面，是一片沙漠區，氣候乾燥，較為貧瘠。古阿拉伯人的居住地靠近古實。古實王謝拉曾率百萬大軍進攻猶大國，而猶大國總共只有五十余萬軍兵。在神的幫助下，亞撒王打敗古實大軍，追趕他們直至基拉耳。
\begin{itemize}
\item 旧约书卷 - 民12:1；代下 12:3；14: 9-15；21:16；王下19:9；
\item 以賽亞、賽20: 2-6；37:9
\item 以西結、29:10；30:4
\item 那鴻等 3:8-10
\item 番2: 12
\item 使徒行传中的埃提阿伯太监
\end{itemize}
旧约许多书卷都論及古實。(民12:1；代下 12:3；14: 9-15；21:16；王下19:9；賽20: 2-6；37:9；結29:10；30:4；番2: 12；鴻3:8-10，等等)案:許多現代譯本己將「古實」一詞譯作「埃塞俄比亞」和「埃塞俄比亞人」。

\subsection{以賽亞書}
\subsubsection{古實之民英勇可畏 18:1-3}
唉！古實河外翅膀刷刷響聲之地， 2 差遣使者在水面上，坐蒲草船過海。先知說：「你們快行的使者，要到高大光滑的民那裡去，自從開國以來那民極其可畏，是分地界踐踏人的，他們的地有江河分開。」 3 世上一切的居民和地上所住的人哪，山上豎立大旗的時候，你們要看！吹角的時候，你們要聽！

\subsubsection{奉禮於耶和華 18:4-7}
耶和華對我這樣說：「我要安靜，在我的居所觀看，如同日光中的清熱，又如露水的雲霧在收割的熱天。」 5 收割之先，花開已謝，花也成了將熟的葡萄，他必用鐮刀削去嫩枝，又砍掉蔓延的枝條。 6 都要撇給山間的鷙鳥和地上的野獸，夏天鷙鳥要宿在其上，冬天野獸都臥在其中。 7 到那時，這高大光滑的民，就是從開國以來極其可畏，分地界踐踏人的，他們的地有江河分開，他們必將禮物奉給萬軍之耶和華，就是奉到錫安山，耶和華安置他名的地方。

\subsubsection{論埃及與古實之預言 20:1-6}
亞述王撒珥根打發他珥探到亞實突的那年，他珥探就攻打亞實突，將城攻取。 2 那時，耶和華曉諭亞摩斯的兒子以賽亞說：「你去，解掉你腰間的麻布，脫下你腳上的鞋。」以賽亞就這樣做，露身赤腳行走。 3 耶和華說：「我僕人以賽亞怎樣露身赤腳行走三年，作為關乎埃及和古實的預兆奇蹟， 4 照樣，亞述王也必擄去埃及人，掠去古實人，無論老少，都露身赤腳，現出下體，使埃及蒙羞。 5 以色列人必因所仰望的古實、所誇耀的埃及，驚惶羞愧。 6 那時，這沿海一帶的居民必說：『看哪，我們素所仰望的，就是我們為脫離亞述王逃往求救的，不過是如此！我們怎能逃脫呢？』」

\subsection{西番雅書 2:12}
古實人哪，你們必被我的刀所殺








\section{埃及之預言}
\subsection{以賽亞書}
\subsubsection{預示埃及之禍亂 19:1-10}
論埃及的默示。

看哪，耶和華乘駕快雲，臨到埃及。埃及的偶像在他面前戰兢，埃及人的心在裡面消化。 2 「我必激動埃及人攻擊埃及人，弟兄攻擊弟兄，鄰舍攻擊鄰舍，這城攻擊那城，這國攻擊那國。 3 埃及人的心神必在裡面耗盡，我必敗壞他們的謀略，他們必求問偶像和念咒的、交鬼的、行巫術的。 4 我必將埃及人交在殘忍主的手中，強暴王必轄制他們。」這是主萬軍之耶和華說的。

5 海中的水必絕盡，河也消沒乾涸。 6 江河要變臭，埃及的河水都必減少枯乾，葦子和蘆荻都必衰殘。 7 靠尼羅河旁的草田，並沿尼羅河所種的田，都必枯乾，莊稼被風吹去，歸於無有。 8 打魚的必哀哭，在尼羅河一切釣魚的必悲傷，在水上撒網的必都衰弱。 9 用梳好的麻造物的，和織白布的，都必羞愧。 10 國柱必被打碎，所有傭工的心必愁煩。

\subsubsection{牧伯謀士盡為愚蒙 19:11-18}
瑣安的首領極其愚昧，法老大有智慧的謀士所籌劃的成為愚謀。你們怎敢對法老說「我是智慧人的子孫，我是古王的後裔」？ 12 你的智慧人在哪裡呢？萬軍之耶和華向埃及所定的旨意，他們可以知道，可以告訴你吧！ 13 瑣安的首領都變為愚昧，挪弗的首領都受了迷惑，當埃及支派房角石的使埃及人走錯了路。 14 耶和華使乖謬的靈摻入埃及中間，首領使埃及一切所做的都有差錯，好像醉酒之人嘔吐的時候東倒西歪一樣。 15 埃及中，無論是頭與尾，棕枝與蘆葦，所做之工都不成就。

16 到那日，埃及人必像婦人一樣，他們必因萬軍之耶和華在埃及以上所掄的手戰兢懼怕。 17 猶大地必使埃及驚恐，向誰提起猶大地，誰就懼怕。這是因萬軍之耶和華向埃及所定的旨意。

18 當那日，埃及地必有五城的人說迦南的方言，又指著萬軍之耶和華起誓，有一城必稱為滅亡城。

\subsubsection{埃及人識主蒙救 19:19-25}
當那日，在埃及地中必有為耶和華築的一座壇，在埃及的邊界上必有為耶和華立的一根柱。 20 這都要在埃及地為萬軍之耶和華做記號和證據。埃及人因為受人的欺壓哀求耶和華，他就差遣一位救主，做護衛者，拯救他們。 21 耶和華必被埃及人所認識，在那日，埃及人必認識耶和華，也要獻祭物和供物敬拜他，並向耶和華許願還願。 22 耶和華必擊打埃及，又擊打又醫治，埃及人就歸向耶和華。他必應允他們的禱告，醫治他們。

23 當那日，必有從埃及通亞述去的大道，亞述人要進入埃及，埃及人也進入亞述，埃及人要與亞述人一同敬拜耶和華。

24 當那日，以色列必與埃及、亞述三國一律，使地上的人得福。 25 因為萬軍之耶和華賜福給他們，說：「埃及我的百姓，亞述我手的工作，以色列我的產業，都有福了！」

\subsubsection{論埃及與古實之預言 20:1-6}
亞述王撒珥根打發他珥探到亞實突的那年，他珥探就攻打亞實突，將城攻取。 2 那時，耶和華曉諭亞摩斯的兒子以賽亞說：「你去，解掉你腰間的麻布，脫下你腳上的鞋。」以賽亞就這樣做，露身赤腳行走。 3 耶和華說：「我僕人以賽亞怎樣露身赤腳行走三年，作為關乎埃及和古實的預兆奇蹟， 4 照樣，亞述王也必擄去埃及人，掠去古實人，無論老少，都露身赤腳，現出下體，使埃及蒙羞。 5 以色列人必因所仰望的古實、所誇耀的埃及，驚惶羞愧。 6 那時，這沿海一帶的居民必說：『看哪，我們素所仰望的，就是我們為脫離亞述王逃往求救的，不過是如此！我們怎能逃脫呢？』」

\subsubsection{戒民倚恃埃及必受羞辱之\textcolor{red}{禍} 30:1-5}
耶和華說：「禍哉，這悖逆的兒女！他們同謀卻不由於我，結盟卻不由於我的靈，以致罪上加罪。 2 起身下埃及去，並沒有求問我，要靠法老的力量加添自己的力量，並投在埃及的蔭下。 3 所以法老的力量必做你們的羞辱，投在埃及的蔭下要為你們的慚愧。 4 他們的首領已在瑣安，他們的使臣到了哈內斯。 5 他們必因那不利於他們的民蒙羞，那民並非幫助，也非利益，只做羞恥凌辱。」


\subsubsection{埃及的幫助是徒然，坐而不動的拉哈伯 30:6-7}
論南方牲畜的默示。
「他們把財物馱在驢駒的脊背上，將寶物馱在駱駝的肉鞍上，經過艱難困苦之地，就是公獅、母獅、蝮蛇、火焰的飛龍之地，往那不利於他們的民那裡去。 7 埃及的幫助是徒然無益的，所以我稱他為『坐而不動的拉哈伯』。

\subsubsection{斥恃埃及不仰望耶和華之\textcolor{red}{禍} 31:1-3}
禍哉，那些下埃及求幫助的！是因仗賴馬匹，倚靠甚多的車輛，並倚靠強壯的馬兵，卻不仰望以色列的聖者，也不求問耶和華。 2 其實耶和華有智慧，他必降災禍，並不反悔自己的話，卻要興起攻擊那作惡之家，又攻擊那作孽幫助人的。 3 埃及人不過是人，並不是神；他們的馬不過是血肉，並不是靈。耶和華一伸手，那幫助人的必絆跌，那受幫助的也必跌倒，都一同滅亡。


\subsection{耶利米 43:1-22}
\subsubsection{預言埃及必為巴比倫王攻擊 43:1-22}
在答比匿，耶和華的話臨到耶利米說： 9 「你在猶大人眼前要用手拿幾塊大石頭藏在砌磚的灰泥中，就是在答比匿法老的宮門那裡， 10 對他們說：『萬軍之耶和華以色列的神如此說：我必召我的僕人巴比倫王尼布甲尼撒來，在所藏的石頭上，我要安置他的寶座，他必將光華的寶帳支搭在其上。 11 他要來攻擊埃及地，定為死亡的必致死亡，定為擄掠的必被擄掠，定為刀殺的必交刀殺。 12 我要在埃及神的廟中使火著起，巴比倫王要將廟宇焚燒，神像擄去。他要得(披上)埃及地，好像牧人披上外衣，從那裡安然而去。 13 他必打碎埃及地伯示麥的柱像，用火焚燒埃及神的廟宇。』」

\subsubsection{責民在埃及崇侍他神必受災禍 44:1-22}
有臨到耶利米的話，論及一切住在埃及地的猶大人，就是住在密奪、答比匿、挪弗、巴忒羅境內的猶大人，說： 2 「萬軍之耶和華以色列的神如此說：我所降於耶路撒冷和猶大各城的一切災禍，你們都看見了。那些城邑今日荒涼，無人居住。 3 這是因居民所行的惡，去燒香侍奉別神，就是他們和你們並你們列祖所不認識的神，惹我發怒。 4 我從早起來差遣我的僕人眾先知去說：『你們切不要行我所厭惡這可憎之事！』 5 他們卻不聽從，不側耳而聽，不轉離惡事，仍向別神燒香。 6 因此我的怒氣和憤怒都倒出來，在猶大城邑中和耶路撒冷的街市上如火著起，以致都荒廢淒涼，正如今日一樣。 7 現在耶和華萬軍之神以色列的神如此說：你們為何做這大惡自害己命，使你們的男人、婦女、嬰孩和吃奶的都從猶大中剪除不留一人呢？ 8 就是因你們手所做的，在所去寄居的埃及地向別神燒香，惹我發怒，使你們被剪除，在天下萬國中令人咒詛、羞辱。 9 你們列祖的惡行，猶大列王和他們后妃的惡行，你們自己和你們妻子的惡行，就是在猶大地、耶路撒冷街上所行的，你們都忘了嗎？ 10 到如今還沒有懊悔，沒有懼怕，沒有遵行我在你們和你們列祖面前所設立的法度、律例。

\subsubsection{預言逃往埃及之猶大人必遭殲滅 44:1-22}
「所以萬軍之耶和華以色列的神如此說：我必向你們變臉降災，以致剪除猶大眾人。 12 那定意進入埃及地在那裡寄居的，就是所剩下的猶大人，我必使他們盡都滅絕。必在埃及地仆倒，必因刀劍、饑荒滅絕，從最小的到至大的都必遭刀劍、饑荒而死，以致令人辱罵、驚駭、咒詛、羞辱。 13 我怎樣用刀劍、饑荒、瘟疫刑罰耶路撒冷，也必照樣刑罰那些住在埃及地的猶大人。 14 甚至那進入埃及地寄居的，就是所剩下的猶大人，都不得逃脫，也不得存留歸回猶大地——他們心中甚想歸回居住之地。除了逃脫的以外，一個都不能歸回。」

\subsubsection{民不聽預言執意仍拜他神 44:1-22}
那些住在埃及地巴忒羅，知道自己妻子向別神燒香的，與旁邊站立的眾婦女，聚集成群，回答耶利米說： 16 「論到你奉耶和華的名向我們所說的話，我們必不聽從！ 17 我們定要成就我們口中所出的一切話，向天后燒香，澆奠祭，按著我們與我們列祖、君王、首領在猶大的城邑中和耶路撒冷的街市上素常所行的一樣。因為那時我們吃飽飯，享福樂，並不見災禍。 18 自從我們停止向天后燒香，澆奠祭，我們倒缺乏一切，又因刀劍、饑荒滅絕。」 19 婦女說：「我們向天后燒香，澆奠祭，做天后像的餅供奉她，向她澆奠祭，是外乎我們的丈夫嗎？」

\subsubsection{耶利米警以必遭諸災 44:1-22}
耶利米對一切那樣回答他的男人、婦女說： 21 「你們與你們列祖、君王、首領並國內的百姓，在猶大城邑中和耶路撒冷街市上所燒的香，耶和華豈不記念，心中豈不思想嗎？ 22 耶和華因你們所做的惡、所行可憎的事，不能再容忍，所以你們的地荒涼，令人驚駭、咒詛，無人居住，正如今日一樣。 23 你們燒香得罪耶和華，沒有聽從他的話，沒有遵行他的律法、條例、法度，所以你們遭遇這災禍，正如今日一樣。」

24 耶利米又對眾民和眾婦女說：「你們在埃及地的一切猶大人，當聽耶和華的話！ 25 萬軍之耶和華以色列的神如此說：你們和你們的妻都口中說、手裡做，說：『我們定要償還所許的願，向天后燒香，澆奠祭。』現在你們只管堅定所許的願而償還吧！ 26 所以你們住在埃及地的一切猶大人，當聽耶和華的話！耶和華說：我指著我的大名起誓：在埃及全地，我的名不再被猶大一個人的口稱呼說『我指著主永生的耶和華起誓』。 27 我向他們留意降禍不降福，在埃及地的一切猶大人必因刀劍、饑荒所滅，直到滅盡。 28 脫離刀劍從埃及地歸回猶大地的，人數很少。那進入埃及地要在那裡寄居的，就是所剩下的猶大人，必知道是誰的話立得住，是我的話呢，是他們的話呢。

\subsubsection{預言埃及見敗於巴比倫為徵 44:1-22}
「耶和華說：我在這地方刑罰你們必有預兆，使你們知道我降禍於你們的話必要立得住。 30 耶和華如此說：我必將埃及王法老合弗拉交在他仇敵和尋索其命的人手中，像我將猶大王西底家交在他仇敵和尋索其命的巴比倫王尼布甲尼撒手中一樣。」

\subsubsection{預言法老必敗於幼發拉底河濱 46:1-12}
耶和華論列國的話臨到先知耶利米。
論到關乎埃及王法老尼哥的軍隊，這軍隊安營在幼發拉底河邊的迦基米施，是巴比倫王尼布甲尼撒在猶大王約西亞的兒子約雅敬第四年所打敗的：

3 「你們要預備大小盾牌，往前上陣！ 4 你們套上車，騎上馬！頂盔站立，磨槍貫甲！ 5 我為何看見他們驚惶轉身退後呢？他們的勇士打敗了，急忙逃跑，並不回頭，驚嚇四圍都有。」這是耶和華說的。 6 「不要容快跑的逃避，不要容勇士逃脫(快跑的不能逃避，勇士不能逃脫)，他們在北方幼發拉底河邊絆跌仆倒。 7 像尼羅河漲發，像江河之水翻騰的，是誰呢？ 8 埃及像尼羅河漲發，像江河的水翻騰，他說：『我要漲發遮蓋遍地，我要毀滅城邑和其中的居民。』 9 馬匹上去吧！車輛急行吧！勇士，就是手拿盾牌的古實人和弗人(呂彼亞人)，並拉弓的路德族，都出去吧！ 10 那日是主萬軍之耶和華報仇的日子，要向敵人報仇。刀劍必吞吃得飽，飲血飲足，因為主萬軍之耶和華在北方幼發拉底河邊有獻祭的事。 11 埃及的民(處女)哪，可以上基列取乳香去！你雖多服良藥，總是徒然，不得治好。 12 列國聽見你的羞辱，遍地滿了你的哀聲，勇士與勇士彼此相碰，一齊跌倒。」

\subsubsection{預言巴比倫王必擊埃及 46:13-26}
耶和華對先知耶利米所說的話，論到巴比倫王尼布甲尼撒要來攻擊埃及地：

14 「你們要傳揚在埃及，宣告在密奪，報告在挪弗、答比匿，說：『要站起出隊，自做準備，因為刀劍在你四圍施行吞滅的事。』 15 你的壯士為何被沖去呢？他們站立不住，因為耶和華驅逐他們。 16 使多人絆跌，他們也彼此撞倒，說：『起來吧！我們再往本民本地去，好躲避欺壓的刀劍。』 17 他們在那裡喊叫說：『埃及王法老不過是個聲音(已經敗亡)，他已錯過所定的時候了。』」 18 君王，名為萬軍之耶和華的說：「我指著我的永生起誓：尼布甲尼撒(他)來的勢派必像他泊在眾山之中，像迦密在海邊一樣。 19 住在埃及的民(女子)哪，要預備擄去時所用的物件！因為挪弗必成為荒場，且被燒毀無人居住。 20 埃及是肥美的母牛犢，但出於北方的毀滅(牛虻)來到了，來到了！ 21 其中的雇勇好像圈裡的肥牛犢，他們轉身退後，一齊逃跑，站立不住，因為他們遭難的日子、追討的時候已經臨到。 22 其中的聲音好像蛇行一樣，敵人要成隊而來，如砍伐樹木的手拿斧子攻擊她。」 23 耶和華說：「埃及的樹林雖然不能尋察(穿不過)，敵人卻要砍伐，因他們多於蝗蟲，不可勝數。 24 埃及的民(女子)必然蒙羞，必交在北方人的手中。」 25 萬軍之耶和華以色列的神說：「我必刑罰挪的亞捫（埃及尊大之神）和法老，並埃及與埃及的神以及君王，也必刑罰法老和倚靠他的人。 26 我要將他們交付尋索其命之人的手和巴比倫王尼布甲尼撒與他臣僕的手，以後埃及必再有人居住，與從前一樣。」這是耶和華說的。

\subsection{以西結書}
\subsubsection{預言埃及王法老遭報 29:1-7}
第十年十月十二日，耶和華的話臨到我說： 2 「人子啊，你要向埃及王法老預言攻擊他和埃及全地， 3 說：『主耶和華如此說：埃及王法老啊，我與你這臥在自己河中的大魚為敵，你曾說：「這河是我的，是我為自己造的。」 4 我耶和華必用鉤子鉤住你的腮頰，又使江河中的魚貼住你的鱗甲。我必將你和所有貼住你鱗甲的魚，從江河中拉上來， 5 把你並江河中的魚都拋在曠野。你必倒在田間，不被收殮，不被掩埋。我已將你給地上野獸、空中飛鳥做食物。 6 埃及一切的居民，因向以色列家成了蘆葦的杖，就知道我是耶和華。 7 他們用手持住你，你就斷折，傷了他們的肩；他們倚靠你，你就斷折，閃了他們的腰。

\subsubsection{必以刀劍滅絕之 29:8-12}
所以主耶和華如此說：我必使刀劍臨到你，從你中間將人與牲畜剪除。 9 埃及地必荒廢淒涼，他們就知道我是耶和華。因為法老說「這河是我的，是我所造的」， 10 所以我必與你並你的江河為敵，使埃及地，從色弗尼塔直到古實境界，全然荒廢淒涼。 11 人的腳、獸的蹄都不經過，四十年之久並無人居住。 12 我必使埃及地在荒涼的國中成為荒涼，使埃及城在荒廢的城中變為荒廢，共有四十年。我必將埃及人分散在列國，四散在列邦。


\subsubsection{越四十年必使微興 29:13-16}
『主耶和華如此說：滿了四十年，我必招聚分散在各國民中的埃及人。 14 我必叫埃及被擄的人回來，使他們歸回本地巴忒羅。在那裡必成為低微的國， 15 必為列國中最低微的，也不再自高於列國之上。我必減少他們，以致不再轄制列國。 16 埃及必不再做以色列家所倚靠的。以色列家仰望埃及人的時候，便思念罪孽，他們就知道我是主耶和華。』

\subsubsection{必以埃及地賜巴比倫王為賞勞 29:17-20}
二十七年正月初一日，耶和華的話臨到我說： 18 「人子啊，巴比倫王尼布甲尼撒使他的軍兵大大效勞，攻打推羅，以致頭都光禿，肩都磨破。然而他和他的軍兵攻打推羅，並沒有從那裡得什麼酬勞。 19 所以主耶和華如此說：我必將埃及地賜給巴比倫王尼布甲尼撒，他必擄掠埃及群眾，搶其中的財為擄物，奪其中的貨為掠物，這就可以做他軍兵的酬勞。 20 我將埃及地賜給他，酬他所效的勞，因王與軍兵是為我勤勞。這是主耶和華說的。

\subsubsection{埃及與其同盟必偕敗亡 30:1-19}
耶和華的話又臨到我說： 2 「人子啊，你要發預言說：『主耶和華如此說：哀哉這日！你們應當哭號。 3 因為耶和華的日子臨近，就是密雲之日，列國受罰之期。 4 必有刀劍臨到埃及，在埃及被殺之人仆倒的時候，古實人就有痛苦；人民必被擄掠，基址必被拆毀。 5 古實人、弗人(呂彼亞人)、路德人、雜族的人民並古巴人，以及同盟之地的人，都要與埃及人一同倒在刀下。

6 『耶和華如此說：扶助埃及的也必傾倒，埃及因勢力而有的驕傲必降低微，其中的人民從色弗尼塔起 b必倒在刀下。這是主耶和華說的。 7 埃及地在荒涼的國中必成為荒涼，埃及城在荒廢的城中也變為荒廢。 8 我在埃及中使火著起，幫助埃及的都被滅絕，那時他們就知道我是耶和華。 9 到那日，必有使者坐船從我面前出去，使安逸無慮的古實人驚懼；必有痛苦臨到他們，好像埃及遭災的日子一樣。看哪，這事臨近了！

10 『主耶和華如此說：我必藉巴比倫王尼布甲尼撒的手除滅埃及眾人。 11 他和隨從他的人，就是列國中強暴的，必進來毀滅這地。他們必拔刀攻擊埃及，使遍地有被殺的人。 12 我必使江河乾涸，將地賣在惡人的手中。我必藉外邦人的手，使這地和其中所有的變為淒涼。這是我耶和華說的。

13 『主耶和華如此說：我必毀滅偶像，從挪弗除滅神像。必不再有君王出自埃及地，我要使埃及地的人懼怕。 14 我必使巴忒羅荒涼，在瑣安中使火著起，向挪施行審判。 15 我必將我的憤怒倒在埃及的保障上，就是訓上，並要剪除挪的眾人。 16 我必在埃及中使火著起，訓必大大痛苦，挪必被攻破，挪弗白日(終日)見仇敵。 17 亞文和比伯實的少年人必倒在刀下，這些城的人必被擄掠。 18 我在答比匿折斷埃及的諸軛，使她因勢力而有的驕傲在其中止息，那時日光必退去。至於這城，必有密雲遮蔽，其中的女子必被擄掠。 19 我必這樣向埃及施行審判，他們就知道我是耶和華。』

\subsubsection{耶和華折法老之臂 30:20-23}
十一年正月初七日，耶和華的話臨到我說： 21 「人子啊，我已打折埃及王法老的膀臂，沒有敷藥，也沒有用布纏好，使他有力持刀。 22 所以主耶和華如此說：看哪，我與埃及王法老為敵，必將他有力的膀臂和已打折的膀臂全行打斷，使刀從他手中墜落。 23 我必將埃及人分散在列國，四散在列邦。

\subsubsection{堅強巴比倫王之臂 30:24-26}
我必使巴比倫王的膀臂有力，將我的刀交在他手中；卻要打斷法老的膀臂，他就在巴比倫王面前唉哼，如同受死傷的人一樣。 25 我必扶持巴比倫王的膀臂，法老的膀臂卻要下垂。我將我的刀交在巴比倫王手中，他必舉刀攻擊埃及地，他們就知道我是耶和華。 26 我必將埃及人分散在列國，四散在列邦，他們就知道我是耶和華。

\subsubsection{向法老述亞述王之榮美 31:1-9}
十一年三月初一日，耶和華的話臨到我說： 2 「人子啊，你要向埃及王法老和他的眾人說：『在威勢上誰能與你相比呢？ 3 亞述王曾如黎巴嫩中的香柏樹，枝條榮美，影密如林，極其高大，樹尖插入雲中。 4 眾水使它生長，深水使它長大，所栽之地有江河圍流，汊出的水道延到田野諸樹。 5 所以它高大超過田野諸樹，發旺的時候枝子繁多，因得大水之力枝條長長。 6 空中的飛鳥都在枝子上搭窩，田野的走獸都在枝條下生子，所有大國的人民都在它蔭下居住。 7 樹大條長，成為榮美，因為根在眾水之旁。 8 神園中的香柏樹不能遮蔽它，松樹不及它的枝子，楓樹不及它的枝條，神園中的樹都沒有它榮美。 9 我使它的枝條蕃多，成為榮美，以致神伊甸園中的樹都嫉妒它。

\subsubsection{因其驕狂邪惡必遭敗亡 31:10-17}
所以主耶和華如此說：因它高大，樹尖插入雲中，心驕氣傲， 11 我就必將它交給列國中大有威勢的人，他必定辦它。我因它的罪惡，已經驅逐它。 12 外邦人，就是列邦中強暴的，將它砍斷棄掉。它的枝條落在山間和一切谷中，它的枝子折斷，落在地的一切河旁。地上的眾民已經走去，離開它的蔭下。 13 空中的飛鳥都要宿在這敗落的樹上，田野的走獸都要臥在它的枝條下， 14 好使水旁的諸樹不因高大而自尊，也不將樹尖插入雲中，並且那些得水滋潤有勢力的也不得高大自立。因為它們在世人中，和下坑的人都被交於死亡，到陰府去了。

15 『主耶和華如此說：它下陰間的那日，我便使人悲哀。我為它遮蓋深淵，使江河凝結，大水停流。我也使黎巴嫩為它悽慘，田野的諸樹都因它發昏。 16 我將它扔到陰間，與下坑的人一同下去，那時列國聽見它墜落的響聲就都震動，並且伊甸的一切樹，就是黎巴嫩得水滋潤最佳最美的樹，都在陰府受了安慰。 17 它們也與它同下陰間，到被殺的人那裡，它們曾做它的膀臂，在列國中它的蔭下居住。

\subsubsection{埃及亦必墮落 31:18}
『在這樣榮耀威勢上，在伊甸園諸樹中，誰能與你相比呢？然而你要與伊甸的諸樹一同下到陰府，在未受割禮的人中與被殺的人一同躺臥。法老和他的群眾乃是如此。這是主耶和華說的。』


\subsubsection{哀嘆埃及之災 32:1-10}
十二年十二月初一日，耶和華的話臨到我說： 2 「人子啊，你要為埃及王法老作哀歌，說：『從前你在列國中如同少壯獅子，現在你卻像海中的大魚。你衝出江河，用爪攪動諸水，使江河渾濁。 3 主耶和華如此說：我必用多國的人民，將我的網撒在你身上，把你拉上來。 4 我必將你丟在地上，拋在田野，使空中的飛鳥都落在你身上，使遍地的野獸吃你得飽。 5 我必將你的肉丟在山間，用你高大的屍首填滿山谷。 6 我又必用你的血澆灌你所游泳之地，漫過山頂，河道都必充滿。 7 我將你撲滅的時候，要把天遮蔽，使眾星昏暗，以密雲遮掩太陽，月亮也不放光。 8 我必使天上的亮光都在你以上變為昏暗，使你的地上黑暗。這是主耶和華說的。 9 我使你敗亡的風聲傳到你所不認識的各國，那時，我必使多民的心因你愁煩。 10 我在許多國民和君王面前向你掄我的刀，國民就必因你驚奇，君王也必因你極其恐慌。在你仆倒的日子，他們各人為自己的性命時刻戰兢。

\subsubsection{必為巴比倫王殲滅 32:11-16}
『主耶和華如此說：巴比倫王的刀必臨到你。 12 我必藉勇士的刀使你的眾民仆倒，這勇士都是列國中強暴的。他們必使埃及的驕傲歸於無有，埃及的眾民必被滅絕。 13 我必從埃及多水旁除滅所有的走獸，人腳、獸蹄必不再攪渾這水。 14 那時，我必使埃及河澄清，江河像油緩流。這是主耶和華說的。 15 我使埃及地變為荒廢淒涼，這地缺少從前所充滿的，又擊殺其中一切的居民，那時他們就知道我是耶和華。 16 人必用這哀歌去哀哭，列國的女子為埃及和他的群眾也必以此悲哀。這是主耶和華說的。』


\subsubsection{埃及之淪亡 32:17-32}
十二年十二月十五日，耶和華的話臨到我說： 18 「人子啊，你要為埃及群眾哀號，又要將埃及和有名之國的女子，並下坑的人一同扔到陰府去。 19 你埃及的美麗勝過誰呢？你下去，與未受割禮的人一同躺臥吧！ 20 他們必在被殺的人中仆倒，她被交給刀劍，要把她和她的群眾拉去。 21 強盛的勇士要在陰間對埃及王和幫助他的說話：『他們是未受割禮被殺的人，已經下去，躺臥不動。』

22 「亞述和她的眾民都在那裡，她民的墳墓在她四圍，他們都是被殺倒在刀下的。 23 他們的墳墓在坑中極深之處，她的眾民在他墳墓的四圍，都是被殺倒在刀下的。他們曾在活人之地使人驚恐。

24 「以攔也在那裡，她的群眾在她墳墓的四圍，都是被殺倒在刀下，未受割禮而下陰府的。他們曾在活人之地使人驚恐，並且與下坑的人一同擔當羞辱。 25 人給她和她的群眾在被殺的人中設立床榻，她民的墳墓在她四圍，他們都是未受割禮被刀殺的。他們曾在活人之地使人驚恐，並且與下坑的人一同擔當羞辱。以攔已經放在被殺的人中。

米設、土巴和她們的群眾都在那裡，她民的墳墓在她四圍，他們都是未受割禮被刀殺的。他們曾在活人之地使人驚恐。 27 他們不得與那未受割禮仆倒的勇士一同躺臥，這些勇士帶著兵器下陰間，頭枕刀劍，骨頭上有本身的罪孽。他們曾在活人之地使勇士驚恐。 28 法老啊，你必在未受割禮的人中敗壞，與那些被殺的人一同躺臥。

29 「以東也在那裡，她君王和一切首領雖然仗著勢力，還是放在被殺的人中。他們必與未受割禮的和下坑的人一同躺臥。 30 在那裡有北方的眾王子和一切西頓人，都與被殺的人下去。他們雖然仗著勢力使人驚恐，還是蒙羞。他們未受割禮，和被刀殺的一同躺臥，與下坑的人一同擔當羞辱。

法老看見他們，便為他被殺的軍隊受安慰。這是主耶和華說的。 32 我任憑法老在活人之地使人驚恐，法老和他的群眾必放在未受割禮和被殺的人中。這是主耶和華說的。

\subsection{预言应验}
\subsubsection{孟裴斯}
首先我们来看看孟裴斯的历史。尤昆哈特（John Urquhart）为我们提供了初步的资料：“埃及人当初命此城名为巴挪弗（Pa-Nouf），巴挪弗也就是今天我们所知的埃及古城孟裴斯。此城据说是米尼士（Menes，注：公元前三十四世纪埃及南北朝的统一者）所建，也是埃及史上第一座开始制定崇拜神像及神庙中敬奉仪式的城市。依此看来，孟裴斯必然是个很受人崇敬的地方。”49/45
　　
　　由孟裴斯早期历史中，我们可以看出这是个以拜偶像闻名的地方。《新旧约全书批评性、试验性与实践性之注释》（A Commentary:Critical,Experimental and Practical on the Old and New Testament;Wm.B.Eerdmans Publishing Co.1961）一书，对孟裴斯曾有如下的记载说：“孟裴斯是埃及中王国（Middle Egypt）的首都，也是‘偶像’之城。历史上虽无明文记载巴比伦王尼布甲尼撒曾毁坏这些偶像，但我们知道后来的波斯王康比西士（Cambyses）在与埃及作战时，曾夺下埃及门户贝鲁辛城（Pelusium）--即圣经中的‘训’城（位于尼罗河上，地中海岸）。他的战略即是将猫狗排在军队之前，因埃及人视猫狗为圣物，因此不肯上前应战，康比西士就这样轻而易举地夺下贝鲁辛城。他入城后屠杀圣牛Apis，并焚毁了许多的偶像。”21/iv 318（参看预言1B，尤哈特把这段史期定在公元前525年，是在先知预言之后）。
\subsubsection{底比斯}
底比斯的历史就完全不同了。根据尤昆哈特的解释，上帝对底比斯的审判极为严厉，历史家们不明白预言，但却很实在地提出预言已被应验的事实。底比斯城曾遭遇空前的大打击，但没有一次是发生在预言之前。尤昆哈特指出以西结生在尼布甲尼撒王在位的时期，在王去世十三年后，波斯人起来统治中东一带，康比西士王在公元前525年进犯埃及，疯狂似地摧毁了底比斯底。他焚毁神庙，打碎高大的偶像，后来底比斯虽然很快地又复兴起来，但已元气大伤，再不能恢复当年的繁荣。49/26，27
　　
　　尤昆哈特又记载底比斯在主耶稣降生前一个世纪时所受的第二次大攻击。虽然经过第一次的毁灭，底比斯仍在那一带诸城当中保持雄厚的经济势力，直到公元前八十九年，敌人围攻这城达三年之久，底比斯才终于倾覆，且从此湮没无闻。该城的土被掘起，正应验了以西结书卅章14、15节的话--底比斯的众人被切铲，再也不得归回本土。49/27

在比较这两座城时我们要记住：底比斯被分裂，城中百姓也被剪除，而孟裴斯城中的偶像要被毁坏，实际上整座城都要被摧毁，韩富乐（Floyd Hamiltom）在《基督徒信仰的根基》（The Basis of Christian Faith）一书在这样解释说：“如果我们比较现在的孟裴斯就会发现，底比斯的偶像依然比比皆是，神庙墙上仍留有许多神像，

\subsubsection{再没有君王从埃及本地出来}
最后一项预言4B常被人所忽略，但是它甚至到今天还可以应用。预言中说再没有君王从埃及本地出来，这点已经全部应验。预言并没有说埃及从此将永远成无政府状态，但是说它的政府要由外族人来领导。果然波斯人在公元前525年征服埃及，埃及在以后的一百七十年间一直是波斯王国下一个判乱的省份。公元前350年，欧克士（Ochus）终于粉碎埃及，从那时起埃及就一直在外国人的统治之下。49/42










\section{阿拉伯的默示}
\subsection{以賽亞書  21:11-12}
論阿拉伯的默示。

底但結伴的客旅啊，你們必在阿拉伯的樹林中住宿。 14 提瑪地的居民拿水來送給口渴的，拿餅來迎接逃避的。 15 因為他們逃避刀劍和出了鞘的刀，並上了弦的弓與刀兵的重災。 16 主對我這樣說：「一年之內，照雇工的年數，基達的一切榮耀必歸於無有， 17 弓箭手所餘剩的，就是基達人的勇士，必然稀少。因為這是耶和華以色列的神說的。」



\section{推羅的默示}
\subsection{以賽亞書}
\subsubsection{預示推羅之荒廢 23:1-11}
論推羅的默示。

他施的船隻，都要哀號！因為推羅變為荒場，甚至沒有房屋，沒有可進之路。這消息是從基提地得來的。 2 沿海的居民，就是素來靠航海西頓的商家得豐盛的，你們當靜默無言！ 3 在大水之上，西曷的糧食、尼羅河的莊稼是推羅的進項，她做列國的大碼頭。 4 西頓哪，你當慚愧！因為大海說，就是海中的保障說：「我沒有劬勞，也沒有生產；沒有養育男子，也沒有撫養童女。」 5 這風聲傳到埃及，埃及人為推羅的風聲極其疼痛。 6 推羅人哪，你們當過到他施去！沿海的居民哪，你們都當哀號！ 7 這是你們歡樂的城，從上古而有的嗎？其中的居民往遠方寄居。

推羅本是賜冠冕的，她的商家是王子，她的買賣人是世上的尊貴人。遭遇如此，是誰定的呢？ 9 是萬軍之耶和華所定的，為要汙辱一切高傲的榮耀，使地上一切的尊貴人被藐視。 10 他施的民 (女)哪，可以流行你的地，好像尼羅河，不再有腰帶拘緊你。 11 耶和華已經向海伸手，震動列國；至於迦南，他已經吩咐拆毀其中的保障。 12 他又說：「受欺壓西頓的居民 (處女)哪，你必不得再歡樂！起來，過到基提去，就是在那裡也不得安歇！」

\subsubsection{越七十年推羅必蒙眷顧 23:12-18}
看哪，迦勒底人之地！向來沒有這民，這國是亞述人為住曠野的人所立的，現在他們建築戍樓，拆毀推羅的宮殿，使她成為荒涼。 14 他施的船隻，都要哀號！因為你們的保障變為荒場。 15 到那時，推羅必被忘記七十年，照著一王的年日。七十年後，推羅的景況必像妓女所唱的歌。 16 你這被忘記的妓女啊，拿琴周流城內！巧彈多唱使人再想念你！ 17 七十年後，耶和華必眷顧推羅，她就仍得利息 (雇價)，與地上的萬國交易 (行淫)。 18 她的貨財和利息要歸耶和華為聖，必不積攢存留，因為她的貨財必為住在耶和華面前的人所得，使他們吃飽，穿耐久的衣服。

\subsection{以西結書}
\subsubsection{推羅狂傲樂耶路撒冷受災亦必遭報 26:1-6}
第十一年十一月初一日，耶和華的話臨到我說： 2 「人子啊，因推羅向耶路撒冷說『啊哈！那做眾民之門的已經破壞，向我開放。她既變為荒場，我必豐盛』， 3 所以主耶和華如此說：推羅啊，我必與你為敵，使許多國民上來攻擊你，如同海使波浪湧上來一樣。 4 他們必破壞推羅的牆垣，拆毀她的城樓。我也要刮淨塵土，使她成為淨光的磐石。 5 她必在海中做曬網的地方，也必成為列國的擄物。這是主耶和華說的。 6 屬推羅城邑的居民 a必被刀劍殺滅，他們就知道我是耶和華。

\subsubsection{推羅將為巴比倫王所攻 26:7-14}
主耶和華如此說：我必使諸王之王的巴比倫王尼布甲尼撒率領馬匹、車輛、馬兵、軍隊和許多人民從北方來攻擊你推羅。 8 他必用刀劍殺滅屬你城邑的居民，也必造臺築壘舉盾牌攻擊你。 9 他必安設撞城錘攻破你的牆垣，用鐵器拆毀你的城樓。 10 因他的馬匹眾多，塵土揚起遮蔽你。他進入你的城門，好像人進入已有破口之城。那時，你的牆垣必因騎馬的和戰車、輜重車的響聲震動。 11 他的馬蹄必踐踏你一切的街道，他必用刀殺戮你的居民，你堅固的柱子 b必倒在地上。 12 人必以你的財寶為擄物，以你的貨財為掠物，破壞你的牆垣，拆毀你華美的房屋，將你的石頭、木頭、塵土都拋在水中。 13 我必使你唱歌的聲音止息，人也不再聽見你彈琴的聲音。 14 我必使你成為淨光的磐石，做曬網的地方，你不得再被建造，因為這是主耶和華說的。

\subsubsection{推羅傾圯海島震動為作哀歌 26:15-21}
主耶和華對推羅如此說：在你中間行殺戮，受傷之人唉哼的時候，因你傾倒的響聲，海島豈不都震動嗎？ 16 那時靠海的君王必都下位，除去朝服，脫下花衣，披上戰兢，坐在地上，時刻發抖，為你驚駭。 17 他們必為你作起哀歌說：『你這有名之城，素為航海之人居住，在海上為最堅固的，平日你和居民使一切住在那裡的人無不驚恐，現在何竟毀滅了！ 18 如今在你這傾覆的日子，海島都必戰兢；海中的群島見你歸於無有，就都驚惶。』

19 「主耶和華如此說：推羅啊，我使你變為荒涼，如無人居住的城邑，又使深水漫過你，大水淹沒你。 20 那時我要叫你下入陰府，與古時的人一同在地的深處，久已荒涼之地居住，使你不再有居民，我也要在活人之地顯榮耀 (在活人之地不再有榮耀)。 21 我必叫你令人驚恐，不再存留於世，人雖尋找你，卻永尋不見。這是主耶和華說的。

\subsubsection{嘆推羅創建完美 27:1-11}
耶和華的話又臨到我說： 2 「人子啊，要為推羅作起哀歌， 3 說：『你居住海口，是眾民的商埠，你的交易通到許多的海島。主耶和華如此說：推羅啊，你曾說：「我是全然美麗的。」 4 你的境界在海中，造你的使你全然美麗。 5 他們用示尼珥的松樹做你的一切板，用黎巴嫩的香柏樹做桅杆， 6 用巴珊的橡樹做你的槳，用象牙鑲嵌基提海島的黃楊木為坐板(艙板)。 7 你的篷帆是用埃及繡花細麻布做的，可以做你的大旗，你的涼篷是用以利沙島的藍色、紫色布做的。 8 西頓和亞發的居民做你盪槳的，推羅啊，你中間的智慧人做掌舵的。 9 迦巴勒的老者和聰明人都在你中間做補縫的，一切泛海的船隻和水手都在你中間經營交易的事。 10 波斯人、路德人、弗人在你軍營中做戰士，他們在你中間懸掛盾牌和頭盔，彰顯你的尊榮。 11 亞發人和你的軍隊都在你四圍的牆上，你的望樓也有勇士，他們懸掛盾牌，成全你的美麗。

\subsubsection{商賈雲集貿易豐富 27:12-25}
他施人因你多有各類的財物，就做你的客商，拿銀、鐵、錫、鉛兌換你的貨物。 13 雅完人、土巴人、米設人都與你交易，他們用人口和銅器兌換你的貨物。 14 陀迦瑪族用馬和戰馬並騾子兌換你的貨物。 15 底但人與你交易，許多海島做你的碼頭，他們拿象牙、烏木與你兌換(進貢)。 16 亞蘭人因你的工作很多，就做你的客商，他們用綠寶石、紫色布繡貨、細麻布、珊瑚、紅寶石兌換你的貨物。 17 猶大和以色列地的人都與你交易，他們用米匿的麥子、餅、蜜、油、乳香兌換你的貨物。 18 大馬士革人因你的工作很多，又因你多有各類的財物，就拿黑本酒和白羊毛與你交易。 19 威但人和雅完人拿紡成的線、亮鐵、桂皮、菖蒲兌換你的貨物。 20 底但人用高貴的毯子、鞍、屜與你交易。 21 阿拉伯人和基達的一切首領都做你的客商，用羊羔、公綿羊、公山羊與你交易。 22 示巴和拉瑪的商人與你交易，他們用各類上好的香料、各類的寶石和黃金兌換你的貨物。 23 哈蘭人、干尼人、伊甸人、示巴的商人和亞述人、基抹人與你交易， 24 這些商人以美好的貨物包在繡花藍色包袱內，又有華麗的衣服裝在香柏木的箱子裡，用繩捆著與你交易。 25 他施的船隻接連成幫為你運貨，你便在海中豐富，極其榮華。

\subsubsection{傾覆淪亡不復振興 27:26-35}
『盪槳的已經把你盪到大水之處，東風在海中將你打破。 27 你的資財、物件、貨物、水手、掌舵的、補縫的、經營交易的，並你中間的戰士和人民，在你破壞的日子必都沉在海中。 28 你掌舵的呼號之聲一發，郊野都必震動。 29 凡盪槳的和水手，並一切泛海掌舵的，都必下船登岸。 30 他們必為你放聲痛哭，把塵土撒在頭上，在灰中打滾。 31 又為你，使頭上光禿，用麻布束腰，號啕痛哭，苦苦悲哀。 32 他們哀號的時候，為你作起哀歌哀哭說：「有何城如推羅？有何城如她在海中成為寂寞的呢？ 33 你由海上運出貨物，就使許多國民充足，你以許多資財、貨物使地上的君王豐富。 34 你在深水中被海浪打破的時候，你的貨物和你中間的一切人民，就都沉下去了。 35 海島的居民為你驚奇，他們的君王都甚恐慌，面帶愁容。 36 各國民中的客商都向你發嘶聲。你令人驚恐，不再存留於世，直到永遠。」』


\subsubsection{推羅君驕侈僭稱神必受重懲 28:15-21}
耶和華的話又臨到我說： 2 「人子啊，你對推羅君王說：『主耶和華如此說：因你心裡高傲，說「我是神，我在海中坐神之位」，你雖然居心自比神，也不過是人，並不是神。 3 （看哪，你比但以理更有智慧，什麼祕事都不能向你隱藏。 4 你靠自己的智慧、聰明得了金銀財寶，收入庫中。 5 你靠自己的大智慧和貿易增添資財，又因資財心裡高傲。） 6 所以主耶和華如此說：因你居心自比神， 7 我必使外邦人，就是列國中的強暴人，臨到你這裡。他們必拔刀砍壞你用智慧得來的美物，褻瀆你的榮光。 8 他們必使你下坑，你必死在海中，與被殺的人一樣。 9 在殺你的人面前，你還能說「我是神」嗎？其實你在殺害你的人手中，不過是人，並不是神。 10 你必死在外邦人手中，與未受割禮(不潔)的人一樣，因為這是主耶和華說的。』

\subsubsection{嘆其榮美為罪汙衊 28:15-21}
耶和華的話臨到我說： 12 「人子啊，你為推羅王作起哀歌，說：『主耶和華如此說：你無所不備，智慧充足，全然美麗。 13 你曾在伊甸神的園中，佩戴各樣寶石，就是紅寶石、紅璧璽、金鋼石、水蒼玉、紅瑪瑙、碧玉、藍寶石、綠寶石、紅玉和黃金，又有精美的鼓、笛在你那裡，都是在你受造之日預備齊全的。 14 你是那受膏遮掩約櫃的基路伯，我將你安置在神的聖山上，你在發光如火的寶石中間往來。 15 你從受造之日所行的都完全，後來在你中間又察出不義。 16 因你貿易很多，就被強暴的事充滿，以致犯罪。所以我因你褻瀆聖地，就從神的山驅逐你。遮掩約櫃的基路伯啊，我已將你從發光如火的寶石中除滅。 17 你因美麗心中高傲，又因榮光敗壞智慧，我已將你摔倒在地，使你倒在君王面前，好叫他們目睹眼見。 18 你因罪孽眾多、貿易不公，就褻瀆你那裡的聖所。故此，我使火從你中間發出，燒滅你，使你在所有觀看的人眼前變為地上的爐灰。 19 各國民中凡認識你的，都必為你驚奇。你令人驚恐，不再存留於世，直到永遠。』


\subsubsection{預言西頓之災 28:20-24}
耶和華的話臨到我說： 21 「人子啊，你要向西頓預言攻擊她， 22 說：『主耶和華如此說：西頓哪，我與你為敵，我必在你中間得榮耀。我在你中間施行審判，顯為聖的時候，人就知道我是耶和華。 23 我必使瘟疫進入西頓，使血流在她街上。被殺的必在其中仆倒，四圍有刀劍臨到她，人就知道我是耶和華。 24 四圍恨惡以色列家的人，必不再向他們做刺人的荊棘、傷人的蒺藜，人就知道我是主耶和華。』

\subsection{阿摩司書 1:9-10}
耶和華如此說：「推羅三番四次地犯罪，我必不免去他的刑罰，因為他將眾民交給以東，並不記念弟兄的盟約。 10 我卻要降火在推羅的城內，燒滅其中的宮殿。」

\subsection{撒迦利亞書 9:1-7}
耶和華的默示應驗在哈得拉地大馬士革——世人和以色列各支派的眼目都仰望耶和華—— 2 和靠近的哈馬，並推羅、西頓，因為這二城的人大有智慧。 3 「推羅為自己修築保障，積蓄銀子如塵沙，堆起精金如街上的泥土。 4 主必趕出她，打敗她海上的權力，她必被火燒滅。 5 亞實基倫看見必懼怕，加沙看見甚痛苦，以革倫因失了盼望蒙羞。加沙必不再有君王，亞實基倫也不再有居民。 6 私生子(外族人)必住在亞實突，我必除滅非利士人的驕傲。 7 我必除去他口中帶血之肉和牙齒內可憎之物，他必作為餘剩的人歸於我們的神，必在猶大像族長，以革倫人必如耶布斯人。


\subsection{推罗的现状-约拿的家}
推罗是以色列北部沿海的一座宏伟的腓尼基城市。但这里的人是崇拜偶像的，他们的王称自己为神，他们以侵扰耶路撒冷为乐。以西结预言了旧推罗城的毁灭：“所以主耶和华如此说：推罗啊，我必与你为敌，使许多国民上来攻击你，如同海使波浪涌上来一样。……我必使你成为净光的磐石，作晒网的地方；你不得再被 建造。因为这是主耶和华说的。”（结26：3，14）

尼布甲尼撒王在大约公元前573年该城毁灭之前，围攻了推罗城十三年之久。在大约公元前525年推罗向波斯人投降。之后在主前332年当亚历山大大帝攻占了离海岸不远的一 个岛后就就完全应验了这一预言。

“亚历山大宣布要视察推罗人的岛屿并向希拉克斯（HERACLES）假神献祭，他们却回答说在推罗大陆的古巴力庙中献祭更合适。当亚历山大坚持必须视察该岛时，他们下令不允许任何波斯人和马其顿人进入他们的城市。由于没有舰队，亚历山大就催毁了旧推罗城，并用其中的瓦砾建造了一条宽200米，横跨海峡的大堤坝，在坝的另一端建起了高塔和战争工事。”

当亚历山大大帝清理旧推罗城，建造他那条出海到岛上的行车道后，旧推罗城变成了像那光秃的岩石。渔民在那里晒干渔网。你可以去看看那一地方，自己亲眼看看圣经预言的逐字应验。在过去的二千三百多年中，旧推罗城还没有重建，正如神所说的判断那样。圣经确实是神所启示的话语。






\section{論以東(度瑪)的默示 }
\subsection{以東（以掃），後裔及以東諸王}
\subsubsection{創世記 25:19-34——以掃雅各生，以掃賣長子之分}
亞伯拉罕的兒子以撒的後代記在下面。亞伯拉罕生以撒。 20 以撒娶利百加為妻的時候正四十歲，利百加是巴旦亞蘭地的亞蘭人彼土利的女兒，是亞蘭人拉班的妹子。 21 以撒因他妻子不生育，就為她祈求耶和華。耶和華應允他的祈求，他的妻子利百加就懷了孕。 22 孩子們在她腹中彼此相爭，她就說：「若是這樣，我為什麼活著呢 c？」她就去求問耶和華。 23 耶和華對她說：「兩國在你腹內，兩族要從你身上出來，這族必強於那族，將來大的要服侍小的。」 24 生產的日子到了，腹中果然是雙子。 25 先產的身體發紅，渾身有毛，如同皮衣，他們就給他起名叫以掃 d。 26 隨後又生了以掃的兄弟，手抓住以掃的腳跟，因此給他起名叫雅各 e。利百加生下兩個兒子的時候，以撒年正六十歲。

兩個孩子漸漸長大，以掃善於打獵，常在田野，雅各為人安靜，常住在帳篷裡。 28 以撒愛以掃，因為常吃他的野味，利百加卻愛雅各。 29 有一天，雅各熬湯，以掃從田野回來累昏了。 30 以掃對雅各說：「我累昏了，求你把這紅湯給我喝。」因此以掃又叫以東 f。 31 雅各說：「你今日把長子的名分賣給我吧。」 32 以掃說：「我將要死，這長子的名分於我有什麼益處呢？」 33 雅各說：「你今日對我起誓吧。」以掃就對他起了誓，把長子的名分賣給雅各。 34 於是雅各將餅和紅豆湯給了以掃，以掃吃了喝了，便起來走了。這就是以掃輕看了他長子的名分。

\subsubsection{創世記 32:19-33:17——雅各以掃相見}
雅各仍舊行路，神的使者遇見他。 2 雅各看見他們就說：「這是神的軍兵。」於是給那地方起名叫瑪哈念 a。 3 雅各打發人先往西珥地去，就是以東地，見他哥哥以掃， 4 吩咐他們說：「你們對我主以掃說：『你的僕人雅各這樣說：我在拉班那裡寄居，直到如今。 5 我有牛、驢、羊群、僕婢，現在打發人來報告我主，為要在你眼前蒙恩。』」 6 所打發的人回到雅各那裡，說：「我們到了你哥哥以掃那裡，他帶著四百人，正迎著你來。」 7 雅各就甚懼怕，而且愁煩，便把那與他同在的人口和羊群、牛群、駱駝分做兩隊， 8 說：「以掃若來擊殺這一隊，剩下的那一隊還可以逃避。」 9 雅各說：「耶和華我祖亞伯拉罕的神，我父親以撒的神啊，你曾對我說：『回你本地本族去，我要厚待你。』 10 你向僕人所施的一切慈愛和誠實，我一點也不配得。我先前只拿著我的杖過這約旦河，如今我卻成了兩隊了。 11 求你救我脫離我哥哥以掃的手，因為我怕他來殺我，連妻子帶兒女一同殺了。 12 你曾說：『我必定厚待你，使你的後裔如同海邊的沙，多得不可勝數。』」
當夜，雅各在那裡住宿，就從他所有的物中拿禮物要送給他哥哥以掃： 14 母山羊二百隻，公山羊二十隻；母綿羊二百隻，公綿羊二十隻； 15 奶崽子的駱駝三十隻，各帶著崽子；母牛四十隻，公牛十隻；母驢二十匹，驢駒十匹。 16 每樣各分一群，交在僕人手下，就對僕人說：「你們要在我前頭過去，使群群相離，有空閒的地方。」 17 又吩咐儘先走的說：「我哥哥以掃遇見你的時候，問你說：『你是哪家的人？要往哪裡去？你前頭這些是誰的？』 18 你就說：『是你僕人雅各的，是送給我主以掃的禮物。他自己也在我們後邊。』」 19 又吩咐第二、第三和一切趕群畜的人說：「你們遇見以掃的時候也要這樣對他說。 20 並且你們要說：『你僕人雅各在我們後邊。』」因雅各心裡說：「我藉著在我前頭去的禮物解他的恨，然後再見他的面，或者他容納我。」 21 於是禮物先過去了，那夜雅各在隊中住宿。
他夜間起來，帶著兩個妻子、兩個使女並十一個兒子，都過了雅博渡口。 23 先打發他們過河，又打發所有的都過去， 24 只剩下雅各一人。有一個人來和他摔跤，直到黎明。 25 那人見自己勝不過他，就將他的大腿窩摸了一把，雅各的大腿窩正在摔跤的時候就扭了。 26 那人說：「天黎明了，容我去吧。」雅各說：「你不給我祝福，我就不容你去。」
那人說：「你名叫什麼？」他說：「我名叫雅各。」 28 那人說：「你的名不要再叫雅各，要叫以色列，因為你與神與人較力，都得了勝。」 29 雅各問他說：「請將你的名告訴我。」那人說：「何必問我的名？」於是在那裡給雅各祝福。 30 雅各便給那地方起名叫毗努伊勒 b，意思說：「我面對面見了神，我的性命仍得保全。」 31 日頭剛出來的時候，雅各經過毗努伊勒，他的大腿就瘸了。 32 故此，以色列人不吃大腿窩的筋，直到今日，因為那人摸了雅各大腿窩的筋。
各舉目觀看，見以掃來了，後頭跟著四百人，他就把孩子們分開交給利亞、拉結和兩個使女， 2 並且叫兩個使女和她們的孩子在前頭，利亞和她的孩子在後頭，拉結和約瑟在儘後頭。 3 他自己在他們前頭過去，一連七次俯伏在地，才就近他哥哥。 4 以掃跑來迎接他，將他抱住，又摟著他的頸項與他親嘴，兩個人就哭了。 5 以掃舉目看見婦人孩子，就說：「這些和你同行的是誰呢？」雅各說：「這些孩子是神施恩給你的僕人的。」 6 於是兩個使女和她們的孩子前來下拜， 7 利亞和她的孩子也前來下拜，隨後約瑟和拉結也前來下拜。 8 以掃說：「我所遇見的這些群畜是什麼意思呢？」雅各說：「是要在我主面前蒙恩的。」 9 以掃說：「兄弟啊，我的已經夠了，你的仍歸你吧。」 10 雅各說：「不然，我若在你眼前蒙恩，就求你從我手裡收下這禮物。因為我見了你的面，如同見了神的面，並且你容納了我。 11 求你收下我帶來給你的禮物，因為神恩待我，使我充足。」雅各再三地求他，他才收下了。 12 以掃說：「我們可以起身前往，我在你前頭走。」 13 雅各對他說：「我主知道孩子們年幼嬌嫩，牛羊也正在乳養的時候，若是催趕一天，群畜都必死了。 14 求我主在僕人前頭走，我要量著在我面前群畜和孩子的力量慢慢地前行，直走到西珥我主那裡。」 15 以掃說：「容我把跟隨我的人留幾個在你這裡。」雅各說：「何必呢？只要在我主眼前蒙恩就是了。」 16 於是，以掃當日起行回往西珥去了。 17 雅各就往疏割去，在那裡為自己蓋造房屋，又為牲畜搭棚，因此那地方名叫疏割 a。

\subsubsection{創世記 36——以掃，西珥之後裔及以東諸王}
以掃，就是以東，他的後代記在下面。 2 以掃娶迦南的女子為妻，就是赫人以倫的女兒亞大和希未人祭便的孫女、亞拿的女兒阿荷利巴瑪， 3 又娶了以實瑪利的女兒、尼拜約的妹子巴實抹。 4 亞大給以掃生了以利法，巴實抹生了流珥， 5 阿荷利巴瑪生了耶烏施、雅蘭、可拉。這都是以掃的兒子，是在迦南地生的。 6 以掃帶著他的妻子、兒女與家中一切的人口，並他的牛羊、牲畜和一切貨財，就是他在迦南地所得的，往別處去，離了他兄弟雅各。 7 因為二人的財物群畜甚多，寄居的地方容不下他們，所以不能同居。 8 於是，以掃住在西珥山裡，以掃就是以東。

9 以掃是西珥山裡以東人的始祖，他的後代記在下面。 10 以掃眾子的名字如下：以掃的妻子亞大生以利法，以掃的妻子巴實抹生流珥。 11 以利法的兒子是提幔、阿抹、洗玻、迦坦、基納斯。 12 亭納是以掃兒子以利法的妾，她給以利法生了亞瑪力。這是以掃的妻子亞大的子孫。 13 流珥的兒子是拿哈、謝拉、沙瑪、米撒。這是以掃妻子巴實抹的子孫。 14 以掃的妻子阿荷利巴瑪是祭便的孫女、亞拿的女兒，她給以掃生了耶烏施、雅蘭、可拉。

以掃子孫中做族長的記在下面。以掃的長子以利法的子孫中，有提幔族長、阿抹族長、洗玻族長、基納斯族長、 16 可拉族長、迦坦族長、亞瑪力族長。這是在以東地從以利法所出的族長，都是亞大的子孫。 17 以掃的兒子流珥的子孫中，有拿哈族長、謝拉族長、沙瑪族長、米撒族長。這是在以東地從流珥所出的族長，都是以掃妻子巴實抹的子孫。 18 以掃的妻子阿荷利巴瑪的子孫中，有耶烏施族長、雅蘭族長、可拉族長。這是從以掃妻子，亞拿的女兒阿荷利巴瑪子孫中所出的族長。 19 以上的族長都是以掃的子孫，以掃就是以東。

那地原有的居民何利人西珥的子孫記在下面，就是羅坍、朔巴、祭便、亞拿、 21 底順、以察、底珊。這是從以東地的何利人西珥子孫中所出的族長。 22 羅坍的兒子是何利、希幔，羅坍的妹子是亭納。 23 朔巴的兒子是亞勒文、瑪拿轄、以巴錄、示玻、阿南。 24 祭便的兒子是亞雅、亞拿（當時在曠野放他父親祭便的驢，遇著溫泉的，就是這亞拿）。 25 亞拿的兒子是底順，亞拿的女兒是阿荷利巴瑪。 26 底順的兒子是欣但、伊是班、益蘭、基蘭。 27 以察的兒子是辟罕、撒番、亞干。 28 底珊的兒子是烏斯、亞蘭。 29 從何利人所出的族長記在下面，就是羅坍族長、朔巴族長、祭便族長、亞拿族長、 30 底順族長、以察族長、底珊族長。這是從何利人所出的族長，都在西珥地，按著宗族做族長。

以色列人未有君王治理以先，在以東地做王的記在下面。 32 比珥的兒子比拉在以東做王，他的京城名叫亭哈巴。 33 比拉死了，波斯拉人謝拉的兒子約巴接續他做王。 34 約巴死了，提幔地的人戶珊接續他做王。 35 戶珊死了，比達的兒子哈達接續他做王，這哈達就是在摩押地殺敗米甸人的，他的京城名叫亞未得。 36 哈達死了，瑪士利加人桑拉接續他做王。 37 桑拉死了，大河邊的利河伯人掃羅接續他做王。 38 掃羅死了，亞革波的兒子巴勒哈南接續他做王。 39 亞革波的兒子巴勒哈南死了，哈達接續他做王，他的京城名叫巴烏，他的妻子名叫米希她別，是米薩合的孫女、瑪特列的女兒。

40 從以掃所出的族長，按著他們的宗族、住處，名字記在下面，就是亭納族長、亞勒瓦族長、耶帖族長、 41 阿何利巴瑪族長、以拉族長、比嫩族長、 42 基納斯族長、提幔族長、米比薩族長、 43 瑪基疊族長、以蘭族長。這是以東人在所得為業的地上，按著他們的住處，所有的族長都是以東人的始祖以掃的後代。

\subsubsection{歷代志上 1:34-54——以掃西珥之後裔，以東諸王族長}
亞伯拉罕生以撒，以撒的兒子是以掃和以色列。以掃的兒子是以利法、流珥、耶烏施、雅蘭、可拉。 36 以利法的兒子是提幔、阿抹、洗玻、迦坦、基納斯、亭納、亞瑪力。 37 流珥的兒子是拿哈、謝拉、沙瑪、米撒。西珥的兒子是羅坍、朔巴、祭便、亞拿、底順、以察、底珊。 39 羅坍的兒子是何利、荷幔，羅坍的妹子是亭納。 40 朔巴的兒子是亞勒文、瑪拿轄、以巴錄、示非、阿南。祭便的兒子是亞雅、亞拿。 41 亞拿的兒子是底順。底順的兒子是哈默蘭、伊是班、益蘭、基蘭。 42 以察的兒子是辟罕、撒番、亞干。底珊的兒子是烏斯、亞蘭。

以色列人未有君王治理之先，在以東地做王的記在下面。有比珥的兒子比拉，他的京城名叫亭哈巴。 44 比拉死了，波斯拉人謝拉的兒子約巴接續他做王。 45 約巴死了，提幔地的人戶珊接續他做王。 46 戶珊死了，比達的兒子哈達接續他做王，這哈達就是在摩押地殺敗米甸人的，他的京城名叫亞未得。 47 哈達死了，瑪士利加人桑拉接續他做王。 48 桑拉死了，大河邊的利河伯人掃羅接續他做王。 49 掃羅死了，亞革波的兒子巴勒哈南接續他做王。 50 巴勒哈南死了，哈達接續他做王，他的京城名叫巴伊，他的妻子名叫米希她別，是米薩合的孫女、瑪特列的女兒。哈達死了，以東人的族長有亭納族長、亞勒瓦族長、耶帖族長、 52 阿何利巴瑪族長、以拉族長、比嫩族長、 53 基納斯族長、提幔族長、米比薩族長、 54 瑪基疊族長、以蘭族長。這都是以東人的族長。

\subsection{出埃及記，民數記 }
\subsubsection{出埃及記 15:13-18——摩西和以色列人向耶和華唱歌以東的族長驚惶}
那時，摩西和以色列人向耶和華唱歌說：「....13 你憑慈愛領了你所贖的百姓，你憑能力引他們到了你的聖所。 14 外邦人聽見就發顫，疼痛抓住非利士的居民。 15 那時，以東的族長驚惶，摩押的英雄被戰兢抓住，迦南的居民心都消化了， 16 驚駭恐懼臨到他們。耶和華啊，因你膀臂的大能，他們如石頭寂然不動，等候你的百姓過去，等候你所贖的百姓過去。 17 你要將他們領進去，栽於你產業的山上。耶和華啊，就是你為自己所造的住處；主啊，就是你手所建立的聖所。 18 耶和華必做王，直到永永遠遠。」

\subsubsection{民數記 20:14-29——亞倫卒在何珥山}
摩西從加低斯差遣使者去見以東王，說：「你的弟兄以色列人這樣說：我們所遭遇的一切艱難， 15 就是我們的列祖下到埃及，我們在埃及久住，埃及人惡待我們的列祖和我們， 16 我們哀求耶和華的時候，他聽了我們的聲音，差遣使者把我們從埃及領出來，這事你都知道。如今，我們在你邊界上的城加低斯。 17 求你容我們從你的地經過。我們不走田間和葡萄園，也不喝井裡的水，只走大道 b，不偏左右，直到過了你的境界。」 18 以東王說：「你不可從我的地經過，免得我帶刀出去攻擊你。」 19 以色列人說：「我們要走大道上去。我們和牲畜若喝你的水，必給你價值。不求別的，只求你容我們步行過去。」 20 以東王說：「你們不可經過。」就率領許多人出來，要用強硬的手攻擊以色列人。 21 這樣，以東王不肯容以色列人從他的境界過去。於是他們轉去，離開他。
\subsubsection{Numbers 21:4Numbers 24:18Numbers 33:37Numbers 34:3Joshua 15:1-25}

\subsection{以賽亞書}
\subsubsection{以賽亞書  21:11-12}
論度瑪的默示。有人聲從西珥呼問我說：「守望的啊，夜裡如何？守望的啊，夜裡如何？」 12 守望的說：「早晨將到，黑夜也來。你們若要問，就可以問，可以回頭再來。」

\subsubsection{耶和華必向萬國復仇申冤 34:1-17}
列國啊，要近前來聽！眾民哪，要側耳而聽！地和其上所充滿的，世界和其中一切所出的，都應當聽！ 2 因為耶和華向萬國發憤恨，向他們的全軍發烈怒，將他們滅盡，交出他們受殺戮。 3 被殺的必然拋棄，屍首臭氣上騰，諸山被他們的血溶化。 4 天上的萬象都要消沒，天被捲起好像書卷；其上的萬象要殘敗，像葡萄樹的葉子殘敗，又像無花果樹的葉子殘敗一樣。 5 「因為我的刀在天上已經喝足，這刀必臨到以東和我所咒詛的民，要施行審判。 6 耶和華的刀滿了血，用脂油和羊羔、公山羊的血，並公綿羊腰子的脂油滋潤的，因為耶和華在波斯拉有獻祭的事，在以東地大行殺戮。 7 野牛、牛犢和公牛要一同下來，他們的地喝醉了血，他們的塵土因脂油肥潤。」

%\subsubsection{以東必世代荒廢 34:8-15}
因耶和華有報仇之日，為錫安的爭辯有報應之年。 9 以東的河水要變為石油，塵埃要變為硫磺，地土成為燒著的石油。 10 晝夜總不熄滅，煙氣永遠上騰，必世世代代成為荒廢，永永遠遠無人經過。 11 鵜鶘、箭豬卻要得為業，貓頭鷹、烏鴉要住在其間。耶和華必將空虛的準繩、混沌的線砣，拉在其上。 12 以東人要召貴胄來治國，那裡卻無一個，首領也都歸於無有。 13 以東的宮殿要長荊棘，保障要長蒺藜和刺草，要做野狗的住處、鴕鳥的居所。 14 曠野的走獸要和豺狼相遇，野山羊要與伴偶對叫，夜間的怪物必在那裡棲身，自找安歇之處。 15 箭蛇要在那裡做窩，下蛋，抱蛋，生子，聚子在其影下，鷂鷹各與伴偶聚集在那裡。你們要查考、宣讀耶和華的書，這都無一缺少，無一沒有伴偶。因為我的口已經吩咐，他的靈將他們聚集。 17 他也為他們拈鬮，又親手用準繩給他們分地，他們必永得為業，世世代代住在其間。

\subsubsection{救主臨格踐踏諸敵 63:1-6}
這從以東的波斯拉來，穿紅衣服，裝扮華美，能力廣大，大步行走的，是誰呢？」「就是我，是憑公義說話，以大能施行拯救。」 2 「你的裝扮為何有紅色，你的衣服為何像踹酒榨的呢？」 3 「我獨自踹酒榨，眾民中無一人與我同在。我發怒將他們踹下，發烈怒將他們踐踏，他們的血濺在我衣服上，並且汙染了我一切的衣裳。 4 因為報仇之日在我心中，救贖我民之年已經來到。 5 我仰望，見無人幫助，我詫異沒有人扶持。所以我自己的膀臂為我施行拯救，我的烈怒將我扶持。 6 我發怒踹下眾民，發烈怒使他們沉醉，又將他們的血倒在地上。

\subsection{耶利米書}
\subsubsection{耶利米書-以東必受懲罰 49:7-22}



\subsubsection{耶利米書-以東必受懲罰 49:7-22}

論以東。萬軍之耶和華如此說：「提幔中再沒有智慧嗎？明哲人不再有謀略嗎？他們的智慧盡歸無有嗎？ 8 底但的居民哪，要轉身逃跑，住在深密處！因為我向以掃追討的時候，必使災殃臨到他。 9 摘葡萄的若來到他那裡，豈不剩下些葡萄呢？盜賊若夜間而來，豈不毀壞直到夠了呢？ 10 我卻使以掃赤露，顯出他的隱密處，他不能自藏。他的後裔、弟兄、鄰舍盡都滅絕，他也歸於無有。 11 你撇下孤兒，我必保全他們的命，你的寡婦可以倚靠我。」 12 耶和華如此說：「原不該喝那杯的一定要喝，你能盡免刑罰嗎？你必不能免，一定要喝。」 13 耶和華說：「我指著自己起誓：波斯拉必令人驚駭、羞辱、咒詛，並且荒涼，她的一切城邑必變為永遠的荒場。」

我從耶和華那裡聽見信息，並有使者被差往列國去說：「你們聚集來攻擊以東，要起來爭戰！」 15 「我使你在列國中為最小，在世人中被藐視。 16 住在山穴中據守山頂的啊，論到你的威嚇，你因心中的狂傲自欺。你雖如大鷹高高搭窩，我卻從那裡拉下你來。」這是耶和華說的。 17 「以東必令人驚駭，凡經過的人就受驚駭，又因她一切的災禍嗤笑。」 18 耶和華說：「必無人住在那裡，也無人在其中寄居，要像所多瑪、蛾摩拉和鄰近的城邑傾覆的時候一樣。 19 仇敵必像獅子從約旦河邊的叢林上來，攻擊堅固的居所。轉眼之間，我要使以東人逃跑，離開這地。誰蒙揀選，我就派誰治理這地。誰能比我呢？誰能給我定規日期呢？有何牧人能在我面前站立得住呢？」

20 你們要聽耶和華攻擊以東所說的謀略和他攻擊提幔居民所定的旨意：仇敵定要將他們群眾微弱的拉去，定要使他們的居所荒涼。 21 因他們仆倒的聲音，地就震動，人在紅海那裡必聽見呼喊的聲音。 22 仇敵必如大鷹飛起，展開翅膀攻擊波斯拉。到那日，以東的勇士心中疼痛，如臨產的婦人。

\subsection{耶利米哀歌-以東必受懲罰 4:17-22}
耶和華發怒成就他所定的，倒出他的烈怒，在錫安使火著起，燒毀錫安的根基。 12 地上的君王和世上的居民都不信敵人和仇敵能進耶路撒冷的城門。 13 這都因她先知的罪惡和祭司的罪孽，他們在城中流了義人的血。 14 他們在街上如瞎子亂走，又被血玷汙，以致人不能摸他們的衣服。 15 人向他們喊著說：「不潔淨的，躲開！躲開！不要挨近我！」他們逃走漂流的時候，列國中有人說：「他們不可仍在這裡寄居。」 16 耶和華發怒，將他們分散，不再眷顧他們。人不重看祭司，也不厚待長老。

17 我們仰望人來幫助，以至眼目失明，還是枉然。我們所盼望的，竟盼望一個不能救人的國。 18 仇敵追趕我們的腳步像打獵的，以致我們不敢在自己的街上行走。我們的結局臨近，我們的日子滿足，我們的結局來到了。 19 追趕我們的比空中的鷹更快，他們在山上追逼我們，在曠野埋伏等候我們。 20 耶和華的受膏者好比我們鼻中的氣，在他們的坑中被捉住，我們曾論到他說：「我們必在他蔭下，在列國中存活。」

\subsection{以西結書}
\subsubsection{摩押，以東，亞捫亦必遭報 25:8-15}
主耶和華如此說：因摩押和西珥人說『看哪，猶大家與列國無異』， 9 所以我要破開摩押邊界上的城邑，就是摩押人看為本國之榮耀的伯耶西末、巴力免、基列亭， 10 好使東方人來攻擊亞捫人。我必將亞捫人之地交給他們為業，使亞捫人在列國中不再被記念。 11 我必向摩押施行審判，他們就知道我是耶和華。

%\subsubsection{以東亦必遭報 25:12-14}
主耶和華如此說：因為以東報仇雪恨，攻擊猶大家，向他們報仇，大大有罪， 13 所以主耶和華如此說：我必伸手攻擊以東，剪除人與牲畜，使以東從提幔起，人必倒在刀下，地要變為荒涼，直到底但。 14 我必藉我民以色列的手報復以東，以色列民必照我的怒氣，按我的憤怒在以東施報，以東人就知道是我施報。這是主耶和華說的。

\subsubsection{預言西珥仇害以色列必遭重罰 35:1-15}
耶和華的話又臨到我說： 2 「人子啊，你要面向西珥山發預言攻擊他， 3 對他說：『主耶和華如此說：西珥山哪，我與你為敵，必向你伸手攻擊你，使你荒涼，令人驚駭。 4 我必使你的城邑變為荒場，成為淒涼，你就知道我是耶和華。 5 因為你永懷仇恨，在以色列人遭災、罪孽到了盡頭的時候，將他們交於刀劍， 6 所以主耶和華說：我指著我的永生起誓，我必使你遭遇流血的報應，罪(血)必追趕你。你既不恨惡殺人流血，所以這罪必追趕你。 7 我必使西珥山荒涼，令人驚駭，來往經過的人我必剪除。 8 我必使西珥山滿有被殺的人，被刀殺的必倒在你小山和山谷，並一切的溪水中。 9 我必使你永遠荒涼，使你的城邑無人居住，你的民就知道我是耶和華。

10 『因為你曾說「這二國這二邦必歸於我，我必得為業」（其實耶和華仍在那裡）， 11 所以主耶和華說：我指著我的永生起誓，我必照你的怒氣和你從仇恨中向他們所發的嫉妒待你。我審判你的時候，必將自己顯明在他們中間。 12 你也必知道我耶和華聽見了你的一切毀謗，就是你攻擊以色列山的話說：『這些山荒涼，是歸我們吞滅的。』 13 你們也用口向我誇大，增添與我反對的話，我都聽見了。 14 主耶和華如此說：全地歡樂的時候，我必使你荒涼。 15 你怎樣因以色列家的地業荒涼而喜樂，我必照你所行的待你。西珥山哪，你和以東全地必都荒涼，你們就知道我是耶和華。』

\subsection{約珥書-以東變為凄凉的曠野 3:19}
埃及必然荒涼，以東變為淒涼的曠野，都因向猶大人所行的強暴，又因在本地流無辜人的血。

\subsection{阿摩司書 1:11-12}
耶和華如此說：「以東三番四次地犯罪，我必不免去他的刑罰，因為他拿刀追趕兄弟，毫無憐憫，發怒撕裂，永懷憤怒。 12 我卻要降火在提幔，燒滅波斯拉的宮殿。」

\subsection{俄巴底亞書 1:1-21}

俄巴底亞得了耶和華的默示。論以東說（我從耶和華那裡聽見信息，並有使者被差往列國去，說：「起來吧，一同起來與以東爭戰」）： 2 「我使你以東在列國中為最小的，被人大大藐視。 3 住在山穴中，居所在高處的啊，你因狂傲自欺，心裡說：『誰能將我拉下地去呢？』 4 你雖如大鷹高飛，在星宿之間搭窩，我必從那裡拉下你來。這是耶和華說的。

5 「盜賊若來在你那裡，或強盜夜間而來（你何竟被剪除！），豈不偷竊直到夠了呢？摘葡萄的若來到你那裡，豈不剩下些葡萄呢？ 6 以掃的隱密處何竟被搜尋？他隱藏的寶物何竟被查出？ 7 與你結盟的都送你上路，直到交界，與你和好的欺騙你，且勝過你，與你一同吃飯的設下網羅陷害你，在你心裡毫無聰明。 8 耶和華說：到那日，我豈不從以東除滅智慧人，從以掃山除滅聰明人？ 9 提幔哪，你的勇士必驚惶，甚至以掃山的人都被殺戮剪除。

10 「因你向兄弟雅各行強暴，羞愧必遮蓋你，你也必永遠斷絕。 11 當外人擄掠雅各的財物，外邦人進入他的城門，為耶路撒冷拈鬮的日子，你竟站在一旁，像與他們同夥。 12 你兄弟遭難的日子，你不當瞪眼看著。猶大人被滅的日子，你不當因此歡樂。他們遭難的日子，你不當說狂傲的話。 13 我民遭災的日子，你不當進他們的城門。他們遭災的日子，你不當瞪眼看著他們受苦。他們遭災的日子，你不當伸手搶他們的財物。 14 你不當站在岔路口，剪除他們中間逃脫的。他們遭難的日子，你不當將他們剩下的人交付仇敵。

15 「耶和華降罰的日子臨近萬國。你怎樣行，他也必照樣向你行，你的報應必歸到你頭上。 16 你們猶大人在我聖山怎樣喝了苦杯，萬國也必照樣常常地喝。且喝且咽，他們就歸於無有。

17 「在錫安山必有逃脫的人，那山也必成聖，雅各家必得原有的產業。 18 雅各家必成為大火，約瑟家必為火焰，以掃家必如碎秸，火必將他燒著吞滅，以掃家必無餘剩的。這是耶和華說的。 19 南地的人必得以掃山，高原的人必得非利士地，也得以法蓮地和撒馬利亞地，便雅憫人必得基列。 20 在迦南人中被擄的以色列眾人必得地直到撒勒法，在西法拉中被擄的耶路撒冷人必得南地的城邑。 21 必有拯救者上到錫安山，審判以掃山，國度就歸耶和華了。」

\subsection{以東人}
\subsubsection{以东的历史}
$https://cclw.net/gospel/new/tzdp/htm/11-2.htm$
乔治·鲁宾逊（Georgl L.Robinson）在其《上古文明之石棺：西拉、以东与以东人》(Sarcophagus of an Ancient Civilization;Petra,Edom And the Ekomites; Macmillian Co.,1930)一书中说到：“当大卫正忙于应付叙利亚北部的琐巴王哈大底谢时，以东常侵犯犹大南部，有时甚至威胁到京都耶路撒冷城。大卫制服以東。有一万八千名的以东人，被杀在死海南端的盛谷 。”44/348

“大卫征服了以东，从此以东便屈服在犹大之下，一直到以色列国分裂，约兰作犹大王时为止。”（参看列王记下八章16至22节）50/179
　　
　　（注：“国际标准圣经百科全书”International Standard Bible Encyclopedia把约兰王在位时期定在公元前八五三年至八四一年，是所有先知在职以前的时期。）20/1580
　　
　　海金写道：“约兰王死后五十年，犹大王亚玛谢（公元前八○○年至八七五年）进犯以东，夺取他们的要塞西拉城。”19/36
　　


\subsubsection{以東人现状-约拿的家}
在巴比伦攻陷耶路撒冷之后，阿拉伯的那巴人（NABATAEAN）攻下了以东。之后以东人生活在犹太地南部。在公元前164年，犹太人击败了他们。在公元前120年犹太人征服了以东族人。大约在公元69年，一支4万人的军队使以东人全部覆灭，从此以东人在历史上消失了，正如先知所预言的那样。

在主后106年罗马征服了以东的 那巴人。他们改变了贸易路线，他们途经PALMYRA 而不走以东首都西拉城了。这就削弱了该地区的经济收入，丧失了重要的经济支柱。BY ZANTINE时期之后，它就变成了荒凉之地直到今日。

在今天你可以去看看这块以东的土地。从北部的提幔直到南部都是荒凉的，正如旧约先知所预言的那样。旅游者可去参观西拉城的遗址。以东成为荒凉之地已有近千年了，这正应验了神的判定。自从新 约使徒时代以来，以东人一个都不存在。应验了旧约先知的预言。证明圣经是神启示的 真道。




\section{耶利米書: 基達必受懲罰 49:28-29}
論巴比倫王尼布甲尼撒所攻打的基達和夏瑣的諸國。耶和華如此說：「迦勒底人哪，起來上基達去，毀滅東方人！ 29 他們的帳篷和羊群都要奪去，將幔子和一切器皿並駱駝為自己掠去。人向他們喊著說：『四圍都有驚嚇！』」



\section{夏瑣Hazor}
\begin{figure}[H] 
\centering
\includegraphics[width=6.in]{ZongJie/DXianZhiFigs/ZongJie/XIASUO}
%\caption{新城地界 1}
\label{fig:sub1}
\end{figure}
\subsection{耶利米書: 夏瑣必受懲罰 49:30-33}
耶和華說：「\textcolor{red}{夏瑣}的居民哪，要逃奔遠方，住在深密處，因為巴比倫王尼布甲尼撒設計謀害你們，起意攻擊你們。」 31 耶和華說：「迦勒底人哪，起來，上安逸無慮的居民那裡去！他們是無門無閂，獨自居住的。 32 他們的駱駝必成為掠物，他們眾多的牲畜必成為擄物。我必將剃周圍頭髮的人分散四方(風)，使災殃從四圍臨到他們。」這是耶和華說的。 33 「\textcolor{red}{夏瑣}必成為野狗的住處，永遠淒涼。必無人住在那裡，也無人在其中寄居。」

\subsection{圣经中的夏瑣其他章节}
\begin{itemize}
\item 書11:1 \textcolor{red}{夏瑣}王耶賓聯合北方諸王，人數多如海邊的沙，來到米倫的水邊，一同安營，要與以色列人爭戰，約書亞在米倫水邊攻打他們，耶和華將他們交在以色列人的手裡，以色列人就擊殺他們，追趕他們直到米斯巴的平原，當時約書亞轉回時，奪了\textcolor{red}{夏瑣}，素來\textcolor{red}{夏瑣}在這諸國中是為首的，以色列人又用火焚燒了\textcolor{red}{夏瑣}。
\item 書12:9 是約書亞所擊殺迦南地的三十一個城邦之王之一。
\item 書19:36 拿弗他利支派所得為業的十九座堅固城之一。
\item 撒上12:9 以色列人又行耶和華眼中看為惡的事，耶和華就把他們付與在\textcolor{red}{夏瑣}作王的迦南王耶賓手中，耶賓王有鐵車九百輛，他大大欺壓以色列人二十年，後來被女先知底波拉所擊敗。
\item 王上9:15 所羅門王挑取服苦的人，建造\textcolor{red}{夏瑣}，米吉多和基色等城。
\item 王下15:29 以色列王比加年間，亞述王提革拉毘列色來奪了以雲，和\textcolor{red}{夏瑣}等城，將居民都擄到亞述去了。
\end{itemize}





\section{以攔}
\begin{figure}[H] 
\centering
\includegraphics[width=6.in]{ZongJie/DXianZhiFigs/ZongJie/ERAM}
%\caption{新城地界 1}
\label{fig:sub1}
\end{figure}
\subsection{耶利米書: 以攔必受懲罰,後日必復振興 25:15-26/49:34-38}
\begin{itemize}
\item 25:15-26 耶和華以色列的神對我如此說：「你從我手中接這杯憤怒的酒，使我所差遣你去的各國的民喝。 16 他們喝了就要東倒西歪，並要發狂，因我使刀劍臨到他們中間。」 17 我就從耶和華的手中接了這杯，給耶和華所差遣我去的各國的民喝，...心利的諸王，\textcolor{red}{以攔}的諸王，瑪代的諸王， 26 北方遠近的諸王以及天下地上的萬國喝了，以後示沙克(巴比倫)王也要喝。
\item 49:34-38 猶大王西底家登基的時候，耶和華論以攔的話臨到先知耶利米說： 35 「萬軍之耶和華如此說：我必折斷以攔人的弓，就是他們為首的權力。 36 我要使四風從天的四方颳來臨到以攔人，將他們分散四方(風)，這被趕散的人沒有一國不到的。 37 耶和華說：我必使以攔人在仇敵和尋索其命的人面前驚惶。我也必使災禍，就是我的烈怒臨到他們。又必使刀劍追殺他們，直到將他們滅盡。 38 我要在以攔設立我的寶座，從那裡除滅君王和首領。這是耶和華說的。「到末後，我還要使被擄的以攔人歸回。這是耶和華說的。」
\end{itemize}

\subsection{結: 以攔被殺倒在刀下 32:22-25}
亞述和她的眾民都在那裡，她民的墳墓在她四圍，他們都是被殺倒在刀下的。 23 他們的墳墓在坑中極深之處，她的眾民在他墳墓的四圍，都是被殺倒在刀下的。他們曾在活人之地使人驚恐。\textcolor{red}{以攔}也在那裡，她的群眾在她墳墓的四圍，都是被殺倒在刀下，未受割禮而下陰府的。他們曾在活人之地使人驚恐，並且與下坑的人一同擔當羞辱。 25 人給她和她的群眾在被殺的人中設立床榻，她民的墳墓在她四圍，他們都是未受割禮被刀殺的。


\subsection{圣经中的以攔其他章节 Elam}
\begin{itemize}
\item 創10:22 閃的兒子是\textcolor{red}{以攔}、亞述、亞法撒、路德、亞蘭。
\item 代上1:17 閃的兒子是\textcolor{red}{以攔}、亞述、亞法撒、路德、亞蘭、烏斯、戶勒、基帖、米設。
\item 創14:1-9 那时，示拿王暗拉非、以拉撒王亚略、\textcolor{red}{以攔}王基大老玛和戈印王提达联合攻打以下五王：所多玛王比拉、蛾摩拉王比沙、押玛王示纳、洗扁王善以别和比拉王，即琐珥王。五王会师西订谷，即盐海。 4 他们受基大老玛王统治十二年，在第十三年叛变了。 5 第十四年，基大老玛联合其他王在亚特律·加宁打败利乏音人，在哈麦打败苏西人，在沙微·基列亭打败以米人， 6 在西珥山打败住在那里的何利人，直杀到靠近旷野的伊勒·巴兰。 7 然后，他们返回安·密巴，即加低斯，征服了亚玛力全境以及住在哈洗逊·他玛的亚摩利人。 8 那时，所多玛王、蛾摩拉王、押玛王、洗扁王和比拉王，即琐珥王，在西订谷摆开阵势， 9 抵抗\textcolor{red}{以攔}王基大老玛、戈印王提达、示拿王暗拉非和以拉撒王亚略：四王跟五王交战。
\item 拉4:9 "省长利宏、书记伸帅和同党的底拿人、亚法萨提迦人、他□拉人、亚法撒人、亚基卫人、巴比伦人、书珊迦人、底亥人、\textcolor{red}{以攔}人,和尊大的亚斯那巴所迁移、安置在撒马利亚城并大河西一带地方的人等，上奏亚达薛西王".
\item 賽11:10-12 到那日，耶西的根立做萬民的大旗，外邦人必尋求他，他安息之所大有榮耀。 11當那日，主必二次伸手救回自己百姓中所餘剩的，就是在亞述、埃及、巴忒羅、古實、\textcolor{red}{以攔}、示拿、哈馬並眾海島所剩下的。 12他必向列國豎立大旗，招回以色列被趕散的人，又從地的四方聚集分散的猶大人。
\item 但8:25 我見了異象的時候，我以為在\textcolor{red}{以攔}省書珊城中，我見異象又如在烏萊河邊。
\item 徒2:5-11 那時，有虔誠的猶太人從天下各國來，住在耶路撒冷。 6 這聲音一響，眾人都來聚集，各人聽見門徒用眾人的鄉談說話，就甚納悶， 7 都驚訝稀奇說：「看哪，這說話的不都是加利利人嗎？ 8 我們各人怎麼聽見他們說我們生來所用的鄉談呢？ 9 我們帕提亞人、瑪代人、\textcolor{red}{以攔}人和住在美索不達米亞、猶太、加帕多家、本都、亞細亞、 10 弗呂家、旁非利亞、埃及的人，並靠近古利奈的利比亞一帶地方的人，從羅馬來的客旅中或是猶太人、或是進猶太教的人， 11 克里特和阿拉伯人，都聽見他們用我們的鄉談，講說神的大作為！」
\end{itemize}










\section{歌革}
\subsection{以西結書38-39}
\subsubsection{預言歌革受災之重 38:1-6}
耶和華的話臨到我說： 2 「人子啊，你要面向瑪各地的歌革，就是羅施、米設、土巴的王發預言攻擊他， 3 說：『主耶和華如此說：羅施、米設、土巴的王歌革啊，我與你為敵。 4 我必用鉤子鉤住你的腮頰，調轉你，將你和你的軍兵、馬匹、馬兵帶出來，都披掛整齊，成了大隊，有大小盾牌，各拿刀劍。 5 波斯人、古實人和弗人(呂彼亞人)各拿盾牌，頭上戴盔， 6 歌篾人和他的軍隊，北方極處的陀迦瑪族和他的軍隊，這許多國的民都同著你。

\subsubsection{述其暴行 38:7-16}
『那聚集到你這裡的各隊都當準備，你自己也要準備，做他們的大帥。 8 過了多日，你必被差派。到末後之年，你必來到脫離刀劍從列國收回之地，到以色列長久荒涼的山上。但那從列國中招聚出來的，必在其上安然居住。 9 你和你的軍隊，並同著你許多國的民，必如暴風上來，如密雲遮蓋地面。

10 『主耶和華如此說：到那時，你心必起意念，圖謀惡計， 11 說：「我要上那無城牆的鄉村，我要到那安靜的民那裡，他們都沒有城牆，無門無閂，安然居住。 12 我去要搶財為擄物，奪貨為掠物，反手攻擊那從前荒涼，現在有人居住之地，又攻擊那住世界中間，從列國招聚，得了牲畜財貨的民。」 13 示巴人、底但人、他施的客商和其間的少壯獅子都必問你說：「你來要搶財為擄物嗎？你聚集軍隊要奪貨為掠物嗎？要奪取金銀，擄去牲畜、財貨嗎？要搶奪許多財寶為擄物嗎？」』

14 「人子啊，你要因此發預言，對歌革說：『主耶和華如此說：到我民以色列安然居住之日，你豈不知道嗎？ 15 你必從本地，從北方的極處率領許多國的民來，都騎著馬，乃一大隊極多的軍兵。 16 歌革啊，你必上來攻擊我的民以色列，如密雲遮蓋地面。末後的日子，我必帶你來攻擊我的地，到我在外邦人眼前，在你身上顯為聖的時候，好叫他們認識我。



\subsubsection{受主烈怒 38:17-23}
『主耶和華如此說：我在古時藉我的僕人以色列的先知所說的，就是你嗎？當日他們多年預言我必帶你來攻擊以色列人。 18 主耶和華說：歌革上來攻擊以色列地的時候，我的怒氣要從鼻孔裡發出。 19 我發憤恨和烈怒如火說，那日在以色列地必有大震動， 20 甚至海中的魚、天空的鳥、田野的獸，並地上的一切昆蟲和其上的眾人，因見我的面就都震動。山嶺必崩裂，陡巖必塌陷，牆垣都必坍倒。 21 主耶和華說：我必命我的諸山發刀劍來攻擊歌革，人都要用刀劍殺害弟兄。 22 我必用瘟疫和流血的事刑罰他，我也必將暴雨、大雹與火並硫磺降於他和他的軍隊，並他所率領的眾民。 23 我必顯為大，顯為聖，在多國人的眼前顯現，他們就知道我是耶和華。』


\subsubsection{預言歌革受罰 39:1-7}
人子啊，你要向歌革發預言攻擊他，說：『主耶和華如此說：羅施、米設、土巴的王歌革啊，我與你為敵。 2 我必調轉你，領你前往，使你從北方的極處上來，帶你到以色列的山上。 3 我必從你左手打落你的弓，從你右手打掉你的箭。 4 你和你的軍隊並同著你的列國人，都必倒在以色列的山上，我必將你給各類的鷙鳥和田野的走獸做食物。 5 你必倒在田野，因為我曾說過。這是主耶和華說的。 6 我要降火在瑪各和海島安然居住的人身上，他們就知道我是耶和華。 7 我要在我民以色列中顯出我的聖名，也不容我的聖名再被褻瀆，列國人就知道我是耶和華，以色列中的聖者。

\subsubsection{以色列人取其器械為薪 39:8-10}
『主耶和華說：這日事情臨近，也必成就，乃是我所說的日子。 9 住以色列城邑的人必出去撿器械，就是大小盾牌、弓箭、梃杖、槍矛，都當柴燒火，直燒七年， 10 甚至他們不必從田野撿柴，也不必從樹林伐木。因為他們要用器械燒火，並且搶奪那搶奪他們的人，擄掠那擄掠他們的人。這是主耶和華說的。


\subsubsection{葬其屍於以色列地 39:11-16}
『當那日，我必將以色列地的谷，就是海東人所經過的谷，賜給歌革為墳地，使經過的人到此停步。在那裡，人必葬埋歌革和他的群眾，就稱那地為哈們歌革谷。 12 以色列家的人必用七個月葬埋他們，為要潔淨全地。 13 全地的居民都必葬埋他們。當我得榮耀的日子，這事必叫他們得名聲。這是主耶和華說的。 14 他們必分派人時常巡查遍地，與過路的人一同葬埋那剩在地面上的屍首，好潔淨全地。過了七個月，他們還要巡查。 15 巡查遍地的人要經過全地，見有人的骸骨，就在旁邊立一標記，等葬埋的人來將骸骨葬在哈們歌革谷。 16 他們必這樣潔淨那地，並有一城名叫哈摩那。』

\subsubsection{飛鳥食其肉飲其血 39:17-20}
人子啊，主耶和華如此說：你要對各類的飛鳥和田野的走獸說：『你們聚集來吧！要從四方聚到我為你們獻祭之地，就是在以色列山上獻大祭之地，好叫你們吃肉喝血。 18 你們必吃勇士的肉，喝地上首領的血，就如吃公綿羊、羊羔、公山羊、公牛，都是巴珊的肥畜。 19 你們吃我為你們所獻的祭，必吃飽了脂油，喝醉了血。 20 你們必在我席上飽吃馬匹和坐車的人，並勇士和一切的戰士。』這是主耶和華說的。




















































 